% Options for packages loaded elsewhere
\PassOptionsToPackage{unicode}{hyperref}
\PassOptionsToPackage{hyphens}{url}
\PassOptionsToPackage{dvipsnames,svgnames,x11names}{xcolor}
%
\documentclass[
  11pt,
]{article}
\usepackage{amsmath,amssymb}
\usepackage{iftex}
\ifPDFTeX
  \usepackage[T1]{fontenc}
  \usepackage[utf8]{inputenc}
  \usepackage{textcomp} % provide euro and other symbols
\else % if luatex or xetex
  \usepackage{unicode-math} % this also loads fontspec
  \defaultfontfeatures{Scale=MatchLowercase}
  \defaultfontfeatures[\rmfamily]{Ligatures=TeX,Scale=1}
\fi
\usepackage{lmodern}
\ifPDFTeX\else
  % xetex/luatex font selection
  \setmainfont[]{Liberation Serif}
  \setsansfont[]{Liberation Sans}
  \setmonofont[]{DejaVu Sans Mono}
\fi
% Use upquote if available, for straight quotes in verbatim environments
\IfFileExists{upquote.sty}{\usepackage{upquote}}{}
\IfFileExists{microtype.sty}{% use microtype if available
  \usepackage[]{microtype}
  \UseMicrotypeSet[protrusion]{basicmath} % disable protrusion for tt fonts
}{}
\makeatletter
\@ifundefined{KOMAClassName}{% if non-KOMA class
  \IfFileExists{parskip.sty}{%
    \usepackage{parskip}
  }{% else
    \setlength{\parindent}{0pt}
    \setlength{\parskip}{6pt plus 2pt minus 1pt}}
}{% if KOMA class
  \KOMAoptions{parskip=half}}
\makeatother
\usepackage{xcolor}
\usepackage[margin=0.8in]{geometry}
\usepackage{longtable,booktabs,array}
\usepackage{calc} % for calculating minipage widths
% Correct order of tables after \paragraph or \subparagraph
\usepackage{etoolbox}
\makeatletter
\patchcmd\longtable{\par}{\if@noskipsec\mbox{}\fi\par}{}{}
\makeatother
% Allow footnotes in longtable head/foot
\IfFileExists{footnotehyper.sty}{\usepackage{footnotehyper}}{\usepackage{footnote}}
\makesavenoteenv{longtable}
\setlength{\emergencystretch}{3em} % prevent overfull lines
\providecommand{\tightlist}{%
  \setlength{\itemsep}{0pt}\setlength{\parskip}{0pt}}
\setcounter{secnumdepth}{-\maxdimen} % remove section numbering
\ifLuaTeX
\usepackage[bidi=basic]{babel}
\else
\usepackage[bidi=default]{babel}
\fi
\babelprovide[main,import]{american}
\ifPDFTeX
\else
\babelfont{rm}[]{Liberation Serif}
\fi
% get rid of language-specific shorthands (see #6817):
\let\LanguageShortHands\languageshorthands
\def\languageshorthands#1{}
\usepackage{adjustbox}
\emergencystretch=3em
\AtBeginEnvironment{longtable}{\small\sloppy}
\usepackage{stmaryrd}
% Printing shims for run-on conditional-bar tokens in the supplied Biology source.
\providecommand{\midL}{\mid L}
\providecommand{\midN}{\mid N}
\sloppy
\ifLuaTeX
  \usepackage{selnolig}  % disable illegal ligatures
\fi
\usepackage{bookmark}
\IfFileExists{xurl.sty}{\usepackage{xurl}}{} % add URL line breaks if available
\urlstyle{same}
\hypersetup{
  pdflang={en-US},
  colorlinks=true,
  linkcolor={blue},
  filecolor={Maroon},
  citecolor={Blue},
  urlcolor={blue},
  pdfcreator={LaTeX via pandoc}}

\author{}
\date{}

\begin{document}

\section{COHERENCE UNIVERSALISM:
PSYCHOLOGY}\label{coherence-universalism-psychology}

\subsection{Meaning, Regulation, and Repair Across
Time}\label{meaning-regulation-and-repair-across-time}

\textbf{Gaura Kiśora Dās Rader}\\
\textbf{Heaven≡Earth Foundation}\\
\textbf{White Paper Series}

\textbf{Version 0.10}\\
\textbf{October 2026}

\textbf{Release status.} This manuscript is part of the complete Version
0.10 proposal. Numbered principles identify the framework's current
formulations; their claim status, external validation, and authorization
for consequential application remain separate questions.

\textbf{Freeze status.} This manuscript is synchronized to the frozen
Version 0.10 release registry, evidence package, and review gates.
Future changes require a new versioned release.

\begin{center}\rule{0.5\linewidth}{0.5pt}\end{center}

\section{Abstract}\label{abstract}

Psychological life is the changing organization of perception, memory,
affect, embodiment, motivation, identity, relationship, action, and
meaning across time. Coherence Universalism describes this domain as
\textbf{meaning-bearing regulation under constraint}. The description
becomes scientifically useful only when the person or other candidate
psychological system, environment, history, protected functions,
observation channels, scale, and temporal horizon are declared.

The paper develops a non-scalar architecture. A psychological model
includes states, environments, possible actions and coping strategies,
transition processes, protected constraints, relational organization,
one or more meaning maps, developmental and relational history, scale,
horizon, and observation model. A meaning map may include explicit
belief, implicit expectation, bodily disposition, affective association,
action affordance, narrative identity, social script, and orienting
commitment. Description-register contents answer to truth and evidence;
commitment-register contents orient identity and action and remain
answerable to fact, integrity, consequence, and revision.

Psychological organization is evaluated through a profile rather than
one universal score. Candidate dimensions include descriptive adequacy,
affect tolerance, regulatory flexibility, identity and narrative
integration, relational capacity, agency, embodied regulation, meaning
and purpose, and repair capacity. These dimensions can conflict. A
worldview may be internally integrated and false. A grief process may
disrupt narrative continuity while preserving reality contact. A coping
strategy may reduce acute distress while narrowing later agency.
Protective compartmentalization may preserve safety and later become
costly. Scalarization therefore requires declared weights, protected
floors, affected-person scope, horizon, and decision purpose; some
profiles should remain incomparable.

The paper treats stress through integration load and regulatory capacity
rather than literal psychological entropy; motivation through plural and
sometimes conflicting aims rather than one coherence drive; trauma
through heterogeneous protective and injurious trajectories rather than
one catastrophe law; neurodiversity through the distinct questions of
difference, disability, distress, danger, mismatch, stigma, and support;
and therapy through plural mechanisms rather than one hidden process.
Self-deception remains central but requires convergent evidence of
motivated asymmetric processing rather than mere error, disagreement,
uncertainty, or defense.

An illustrative model demonstrates a horizon-sensitive reversal between
immediate suppression and supported engagement under declared
assumptions. It is a pedagogical construction, not clinical evidence.
The broader empirical program asks whether profile-based and
history-sensitive models improve prediction, intervention, and
conceptual clarity beyond the constructs they translate.

The paper closes by locating healing within the interpretive and
commitment registers. Psychological healing is sacred work when it
restores truthful contact, embodied capacity, agency, relationship, and
the ability to hold experience without being wholly governed by what
once had to be survived. This interpretation never sanctifies trauma or
turns spiritual meaning into clinical authority.

\textbf{Keywords:} psychology; meaning; regulation; trauma;
self-deception; attachment; neurodiversity; psychotherapy; repair;
sacred healing

\subsection{Relation to the series}\label{relation-to-the-series}

Biology provides active self-maintenance and organism-relative stake.
Consciousness distinguishes access, perspective, valence, temporal
depth, identity, and ego. Psychology begins where a conscious or
behaviorally organized being forms learned, embodied, socially shaped,
action-guiding relations among events, self, others, values, and
futures. Social Dynamics and Ethics extend these processes into
institutions, power, justification, and repair.

This paper can be read independently. The
\href{../supplements/06_Psychology_Technical_Supplement_v0.10.md}{technical
supplement} preserves the complete formal propositions, claim ledger,
candidate-principle registry, deterministic model, independent audit,
and revision history. This paper is not a clinical manual: it does not
diagnose individuals or recommend treatment for particular persons.

\subsection{Note on Version 0.10}\label{note-on-version-0.10}

Version 0.10 retains the original concern with integration, meaning,
trauma, attachment, self-deception, and healing while replacing
universal scores, gradients, entropy identities, developmental
hierarchies, and clinical overreach with profiles, declared models,
protected floors, plural mechanisms, and explicit register discipline.
The main text presents the current account directly; Appendix B and the
technical archive preserve the revision record.\footnote{The historical
  archive preserves the v0.09 paper, systematic reconstruction,
  independent audit, audit response, and complete migration record. The
  publication manuscript does not require those materials for ordinary
  comprehension.}

\textbf{Release-candidate synchronization.} This manuscript uses the
shared Version 0.10 definitions, principle identifiers, Mandala
cross-references, and interpretive and commitment language published in
\emph{Coherence Universalism: Metaphysics and Epistemology} and
\emph{Coherence Universalism: Foundations}. Its audited technical core,
formal results, numerical results, and claim-status boundaries remain
unchanged except where this manuscript explicitly states otherwise.

\subsection{How claims are marked}\label{how-claims-are-marked}

The paper uses \textbf{definition}, \textbf{formal result},
\textbf{imported result}, \textbf{empirical generalization},
\textbf{conditional result}, \textbf{bridge principle},
\textbf{hypothesis}, \textbf{interpretation}, \textbf{commitment}, and
\textbf{open problem} in their ordinary scoped senses. A domain
authority identifies the controlling CU formulation; it does not convert
a conditional or provisional claim into established fact.

\section{Contents}\label{contents}

\begin{itemize}
\tightlist
\item
  \begin{enumerate}
  \def\labelenumi{\arabic{enumi}.}
  \tightlist
  \item
    Psychology as Meaning-Bearing Regulation
  \end{enumerate}
\item
  \begin{enumerate}
  \def\labelenumi{\arabic{enumi}.}
  \setcounter{enumi}{1}
  \tightlist
  \item
    A Non-Scalar Architecture for Psychological Organization
  \end{enumerate}
\item
  \begin{enumerate}
  \def\labelenumi{\arabic{enumi}.}
  \setcounter{enumi}{2}
  \tightlist
  \item
    Perception, Memory, and the Construction of Meaning
  \end{enumerate}
\item
  \begin{enumerate}
  \def\labelenumi{\arabic{enumi}.}
  \setcounter{enumi}{3}
  \tightlist
  \item
    Affect, Motivation, Embodiment, and Agency
  \end{enumerate}
\item
  \begin{enumerate}
  \def\labelenumi{\arabic{enumi}.}
  \setcounter{enumi}{4}
  \tightlist
  \item
    Identity, Personality, and Self-Deception
  \end{enumerate}
\item
  \begin{enumerate}
  \def\labelenumi{\arabic{enumi}.}
  \setcounter{enumi}{5}
  \tightlist
  \item
    Attachment, Sexuality, Intimacy, and Co-Regulation
  \end{enumerate}
\item
  \begin{enumerate}
  \def\labelenumi{\arabic{enumi}.}
  \setcounter{enumi}{6}
  \tightlist
  \item
    Stress, Trauma, Dissociation, and Resilience
  \end{enumerate}
\item
  \begin{enumerate}
  \def\labelenumi{\arabic{enumi}.}
  \setcounter{enumi}{7}
  \tightlist
  \item
    Neurodiversity, Psychopathology, and Diagnosis
  \end{enumerate}
\item
  \begin{enumerate}
  \def\labelenumi{\arabic{enumi}.}
  \setcounter{enumi}{8}
  \tightlist
  \item
    Therapy, Medication, and Other Interventions
  \end{enumerate}
\item
  \begin{enumerate}
  \def\labelenumi{\arabic{enumi}.}
  \setcounter{enumi}{9}
  \tightlist
  \item
    Development, Plasticity, and Psychological Trajectories
  \end{enumerate}
\item
  \begin{enumerate}
  \def\labelenumi{\arabic{enumi}.}
  \setcounter{enumi}{10}
  \tightlist
  \item
    Psychological Infrastructure and the Meaning-Crisis Hypothesis
  \end{enumerate}
\item
  \begin{enumerate}
  \def\labelenumi{\arabic{enumi}.}
  \setcounter{enumi}{11}
  \tightlist
  \item
    Empirical Program, Measurement, and Falsification
  \end{enumerate}
\item
  \begin{enumerate}
  \def\labelenumi{\arabic{enumi}.}
  \setcounter{enumi}{12}
  \tightlist
  \item
    Psychology, Healing, and Sacred Repair
  \end{enumerate}
\item
  \begin{enumerate}
  \def\labelenumi{\arabic{enumi}.}
  \setcounter{enumi}{13}
  \tightlist
  \item
    What Commitments This Account Licenses
  \end{enumerate}
\item
  \begin{enumerate}
  \def\labelenumi{\arabic{enumi}.}
  \setcounter{enumi}{14}
  \tightlist
  \item
    Conclusion: Truthful Life Across Change
  \end{enumerate}
\item
  Appendix A. Psychology Principles and Claim-Governance Summary
\item
  Appendix B. Major Revisions from v0.09
\item
  Appendix C. Technical Supplement and Reproducibility
\item
  References
\end{itemize}

\section{Principal Contributions}\label{principal-contributions}

\begin{enumerate}
\def\labelenumi{\arabic{enumi}.}
\tightlist
\item
  Defines psychological organization as meaning-bearing regulation under
  declared constraints.
\item
  Replaces a universal psychological scalar with a multidimensional
  profile and partial-order discipline.
\item
  Separates descriptive accuracy, existential meaning, commitment,
  affect, functional stake, and phenomenal valence.
\item
  Treats protective compartmentalization before judging fragmentation or
  pathology.
\item
  Reconstructs trauma, dissociation, resilience, and recovery as
  heterogeneous and history-dependent.
\item
  Preserves self-deception as a central process while tightening its
  evidential standard.
\item
  Separates therapy outcomes and mechanisms and limits CU's clinical
  authority.
\item
  Develops healing as sacred repair without sanctifying trauma.
\end{enumerate}

\begin{center}\rule{0.5\linewidth}{0.5pt}\end{center}

\section{Domain Analytic Frame and CAP
Record}\label{domain-analytic-frame-and-cap-record}

\textbf{Target:} Coherence Universalism Version 0.10 release\\
\textbf{Status:} Release declaration; scope-limited and revisable under
the declared review gates\\
\textbf{Use rule:} Unresolved fields are left explicit rather than
completed by analogy.

\subsection{Paper and model identity}\label{paper-and-model-identity}

\begin{longtable}[]{@{}
  >{\raggedright\arraybackslash}p{(\columnwidth - 4\tabcolsep) * \real{0.3333}}
  >{\raggedright\arraybackslash}p{(\columnwidth - 4\tabcolsep) * \real{0.3333}}
  >{\raggedright\arraybackslash}p{(\columnwidth - 4\tabcolsep) * \real{0.3333}}@{}}
\toprule\noalign{}
\begin{minipage}[b]{\linewidth}\raggedright
Field
\end{minipage} & \begin{minipage}[b]{\linewidth}\raggedright
Declaration
\end{minipage} & \begin{minipage}[b]{\linewidth}\raggedright
Status
\end{minipage} \\
\midrule\noalign{}
\endhead
\bottomrule\noalign{}
\endlastfoot
Paper & Coherence Universalism: Psychology & Specified \\
Domain owner & Psychology & Specified by the Version 0.10 release
corpus \\
Primary model or analysis & Psychological model
\(\mathcal P=(X,E,U,F,K,G,\Phi,H,\pi,T,\mathcal O)\) & Paper-relative \\
Primary claim status & Definitions, empirical generalizations, bridge
principles, hypotheses, illustrative models, interpretations, and
commitments & Mixed statuses remain explicit \\
Review gate & Psychology, psychiatry, psychotherapy, neuroscience,
measurement, disability and neurodiversity review, and affected-person
participation & Proportionate to claim and consequential reach \\
Methodological authority & CU-FP2, CU-M2, CU-M11, CU-M12, CU-M13,
CU-M14, and CU-M15 & Controlling Version 0.10 methodological
framework \\
\end{longtable}

\subsection{Shared Analytic Frame}\label{shared-analytic-frame}

\begin{longtable}[]{@{}
  >{\raggedright\arraybackslash}p{(\columnwidth - 4\tabcolsep) * \real{0.3333}}
  >{\raggedright\arraybackslash}p{(\columnwidth - 4\tabcolsep) * \real{0.3333}}
  >{\raggedright\arraybackslash}p{(\columnwidth - 4\tabcolsep) * \real{0.3333}}@{}}
\toprule\noalign{}
\begin{minipage}[b]{\linewidth}\raggedright
\#
\end{minipage} & \begin{minipage}[b]{\linewidth}\raggedright
Slot and domain declaration
\end{minipage} & \begin{minipage}[b]{\linewidth}\raggedright
Status and evidence
\end{minipage} \\
\midrule\noalign{}
\endhead
\bottomrule\noalign{}
\endlastfoot
1 & \textbf{Unit and boundary.} A person or other candidate
psychological system embedded in bodily, relational, cultural, material,
and institutional environments. & Must be declared \\
2 & \textbf{State or trajectory space.} \(X\), meaning maps \(\Phi\),
and history \(H\) represent changing regulation, identity, relation, and
action across time. & Model-specific \\
3 & \textbf{Environment.} \(E\) includes physical, interpersonal,
cultural, economic, and institutional conditions rather than treating
them as inner variables. & Explicit \\
4 & \textbf{Constraints and viable region.} \(K\) includes safety,
consent, agency, bodily function, relational, and other protected
conditions; distress, difference, disability, and disorder remain
distinct. & Plural \\
5 & \textbf{Actions, options, and transitions.} \(U\) includes coping,
action, support, refusal, treatment, environmental change, and
developmental transitions. & Options can be structurally constrained \\
6 & \textbf{Dynamics or update rule.} \(F\) may be qualitative,
stochastic, developmental, relational, or data-driven; no universal
psychological flow is asserted. & Model-specific \\
7 & \textbf{Observation and access.} \(\mathcal O\) links report,
behavior, physiology, context, history, and third-person assessment; no
channel is complete. & Multimethod \\
8 & \textbf{Perturbation and intervention class.} Stressors, relational
events, environmental changes, medication, therapy, exposure, support,
movement, and policy can be interventions. & Purpose- and
safety-bound \\
9 & \textbf{Equilibria, basins, or regimes.} Attractor language may
apply to particular models, but the paper supplies no general
psychological phase portrait or universal basin evidence. & Not
generally specified \\
10 & \textbf{Local--global status.} Person, relationship, group,
institution, and horizon can reverse judgments; profile coordinates are
compared separately unless licensed. & Qualitative / coordinate-wise \\
11 & \textbf{Scalar status.} The primary object is a nine-coordinate
profile. Scores are measurement or decision instruments for declared
purposes only. & No scalar by default \\
12 & \textbf{Statistical-image and allocation status.} Belief
distributions, time use, attention shares, and some occupancy models may
qualify. Psychological profiles and meaning maps do not become
distributions through normalization. & Candidate subobjects only \\
13 & \textbf{Translation status.} Allostatic, control, predictive,
narrative, clinical, contemplative, and social constructs retain their
own assumptions and evidence. & Levels 1--4 by construct \\
14 & \textbf{Open gaps.} Profile validation, within-person dynamics,
cultural transportability, treatment mechanisms, horizon effects, and
structural causation remain active research problems. & Open \\
\end{longtable}

\subsection{Compact Coherence Analysis Protocol
record}\label{compact-coherence-analysis-protocol-record}

\begin{longtable}[]{@{}
  >{\raggedright\arraybackslash}p{(\columnwidth - 4\tabcolsep) * \real{0.3333}}
  >{\raggedright\arraybackslash}p{(\columnwidth - 4\tabcolsep) * \real{0.3333}}
  >{\raggedright\arraybackslash}p{(\columnwidth - 4\tabcolsep) * \real{0.3333}}@{}}
\toprule\noalign{}
\begin{minipage}[b]{\linewidth}\raggedright
CAP module
\end{minipage} & \begin{minipage}[b]{\linewidth}\raggedright
Declaration
\end{minipage} & \begin{minipage}[b]{\linewidth}\raggedright
Boundary
\end{minipage} \\
\midrule\noalign{}
\endhead
\bottomrule\noalign{}
\endlastfoot
System declaration & State person or unit, environment, history, meaning
maps, protected functions, scale, horizon, actions, and observation
model. & Section 2.1 \\
Observation declaration & Separate self-report, behavior, physiology,
clinician judgment, relational testimony, and material context. & No one
privileged channel \\
Evaluation declaration & Declare profile dimensions, floors,
uncertainty, affected-person scope, horizon, scalar purpose, and
non-dominance. & Screen before score \\
Power and authorship declaration & Name who defines health, impairment,
treatment goals, diagnosis, evidence, and acceptable risk; preserve
consent and refusal. & Affected-person standing required \\
Revision declaration & Changed goals, environment, diagnosis, floor,
horizon, or identity criterion constitute model revision rather than
mere new data. & Version and re-evaluate \\
\end{longtable}

\subsection{License summary}\label{license-summary}

\begin{longtable}[]{@{}
  >{\raggedright\arraybackslash}p{(\columnwidth - 4\tabcolsep) * \real{0.3333}}
  >{\raggedright\arraybackslash}p{(\columnwidth - 4\tabcolsep) * \real{0.3333}}
  >{\raggedright\arraybackslash}p{(\columnwidth - 4\tabcolsep) * \real{0.3333}}@{}}
\toprule\noalign{}
\begin{minipage}[b]{\linewidth}\raggedright
License
\end{minipage} & \begin{minipage}[b]{\linewidth}\raggedright
Declaration
\end{minipage} & \begin{minipage}[b]{\linewidth}\raggedright
Current status
\end{minipage} \\
\midrule\noalign{}
\endhead
\bottomrule\noalign{}
\endlastfoot
Scalar license & Measurement scalars for specific constructs and
purpose-bound decision scores only. & Conditional \\
Local--global license & Coordinate-wise or qualitative; group averages
do not establish within-person dynamics. & Default \\
Statistical Image License & Possible for declared beliefs, predictions,
choices, time-use distributions, or probabilistic response models. &
Object-specific \\
Allocation License & Possible for real time or attention budgets;
regulatory capacity is not assumed to be a fixed divisible total. &
Conditional empirical claim \\
Refinement and Coarse-Graining License & Permitted only when categories
and their psychological meaning remain stable enough under aggregation.
& Conditional \\
Fisher-Rao License & Possible for licensed statistical submodels, not
for the nine-coordinate psychological profile or meaning space as a
whole. & Object-specific only \\
Translation license & Imported clinical constructs require
target-population validation; contemplative interpretations do not
inherit empirical status. & Mixed \\
\end{longtable}

\subsection{Open gaps and review
disposition}\label{open-gaps-and-review-disposition}

\begin{longtable}[]{@{}
  >{\raggedright\arraybackslash}p{(\columnwidth - 4\tabcolsep) * \real{0.3333}}
  >{\raggedright\arraybackslash}p{(\columnwidth - 4\tabcolsep) * \real{0.3333}}
  >{\raggedright\arraybackslash}p{(\columnwidth - 4\tabcolsep) * \real{0.3333}}@{}}
\toprule\noalign{}
\begin{minipage}[b]{\linewidth}\raggedright
Gap
\end{minipage} & \begin{minipage}[b]{\linewidth}\raggedright
Consequence
\end{minipage} & \begin{minipage}[b]{\linewidth}\raggedright
Required work
\end{minipage} \\
\midrule\noalign{}
\endhead
\bottomrule\noalign{}
\endlastfoot
Clinical validation & The framework is not a diagnostic or treatment
manual. & Prospective and comparative studies \\
Cultural and disability scope & Unexamined norms can turn difference
into pathology. & Participatory and cross-cultural review \\
Structural causation & Individual profiles can conceal social production
of distress. & Link person-level and institutional evidence without
reduction \\
\end{longtable}

\begin{quote}
\textbf{Propagation rule.} This card makes the paper's existing objects
and limits visible. It does not strengthen any empirical, mathematical,
normative, metaphysical, or spiritual claim beyond the wording supported
in the manuscript.
\end{quote}

\begin{center}\rule{0.5\linewidth}{0.5pt}\end{center}

\section{1. Psychology as meaning-bearing
regulation}\label{psychology-as-meaning-bearing-regulation}

\subsection{1.1 The integrative problem}\label{the-integrative-problem}

Psychology studies processes that are separated by method but joined in
lived experience. Perception changes what can be remembered. Memory
changes what a situation means. Meaning changes affect and action.
Attachment changes the cost of uncertainty. Identity changes which
evidence is threatening. Material conditions change the regulatory
options available. Culture supplies categories through which distress
becomes intelligible or invisible.

The field's pluralism is not merely a defect. Neural, behavioral,
phenomenological, interpersonal, developmental, cultural, and
institutional questions require different methods. The problem begins
when one useful construct is treated as the hidden essence of all the
others.

Coherence is proposed here as a comparative vocabulary for relations
among processes, not as one psychological substance or score.

\subsection{1.2 What the psychological domain
adds}\label{what-the-psychological-domain-adds}

Biology supplies active self-maintenance and organism-relative stake.
Consciousness supplies candidate perspective, temporal organization, and
possibly valence. Psychology begins where an experiencing or
behaviorally organized system forms learned, socially shaped,
action-guiding relations among events, self, others, possibilities,
values, and futures.

\textbf{{[}Definition{]} Psychological organization} is the temporally
extended coordination of perception, affect, memory, embodiment,
motivation, self- and other-modeling, relationship, and action in a
person or other candidate psychological system.

\textbf{{[}Definition{]} Meaning-bearing regulation} is regulation in
which states, events, or possibilities function not only as physical
causes but as signs of relevance, affordance, threat, opportunity,
obligation, identity, belonging, or anticipated consequence.

Meaning includes semantic content, situational implication, motivational
significance, narrative place, existential orientation, and
commitment-bearing practice. These uses interact without becoming
identical.

\subsection{1.3 The bounded thesis}\label{the-bounded-thesis}

\begin{quote}
\textbf{{[}Bridge principle{]} Psychological explanation often improves
when processes are studied as parts of a constrained, history-dependent
organization whose components can support, conflict with, shield, or
destabilize one another across scales and horizons.}
\end{quote}

The claim earns content only where the organizational description
improves prediction, intervention, or conceptual clarity beyond existing
constructs. It should be narrowed or abandoned when coherence language
merely renames what is already known, cannot discriminate among
competing explanations, or repeatedly licenses judgments insensitive to
culture, disability, power, and individual aims.

\subsection{1.4 Conceptual prohibitions}\label{conceptual-prohibitions}

Every later claim is constrained by the following rules:

\begin{enumerate}
\def\labelenumi{\arabic{enumi}.}
\tightlist
\item
  \textbf{No scalar by default.} A profile may be collapsed only with
  declared weights, protected floors, uncertainty, affected-person
  scope, and decision purpose.
\item
  \textbf{No health from consistency alone.} Internal consistency can
  stabilize delusion, prejudice, compulsion, abuse, or despair.
\item
  \textbf{No pathology from difference alone.} Statistical atypicality
  and social nonconformity are not equivalent to disorder.
\item
  \textbf{No self-deception from disagreement alone.} Falsehood,
  defense, uncertainty, and motivated distortion require different
  evidence.
\item
  \textbf{No individualization of structural harm.} Poverty,
  discrimination, violence, exclusion, and institutional betrayal cannot
  be understood solely as inner deficit.
\item
  \textbf{No spiritual bypass.} Meaning, acceptance, oneness, or
  nonattachment may not erase grief, anger, bodily limits, injustice, or
  boundaries.
\item
  \textbf{No clinical authority from framework language.} CU does not
  replace diagnosis, informed consent, professional competence, or
  established safety standards.
\end{enumerate}

\subsection{1.5 Roadmap}\label{roadmap}

Sections 2--6 develop the profile, meaning, affect, identity, and
relational architecture. Sections 7--10 address trauma, diagnosis,
intervention, and development. Section 11 situates minds within social
infrastructure. Section 12 gives the empirical program. Sections 13 and
14 develop the interpretive and commitment registers. Section 15 gathers
the scientific and spiritual account into one bounded conclusion.

\begin{center}\rule{0.5\linewidth}{0.5pt}\end{center}

\section{2. A non-scalar architecture for psychological
organization}\label{a-non-scalar-architecture-for-psychological-organization}

\subsection{2.1 Declaring a psychological
model}\label{declaring-a-psychological-model}

A candidate psychological model is represented as

\[\adjustbox{max width=\linewidth}{$\displaystyle 
\mathcal P
=
(X,E,U,F,K,G,\Phi,H,\pi,T,\mathcal O).
$}\]

The entries are:

\begin{itemize}
\tightlist
\item
  \(X\): state variables or latent processes included in the model;
\item
  \(E\): relevant physical, interpersonal, cultural, and institutional
  environment;
\item
  \(U\): available actions, coping strategies, or interventions;
\item
  \(F\): transition dynamics, which may be deterministic, stochastic,
  qualitative, or data-driven;
\item
  \(K\): protected constraints, safety floors, or admissibility
  conditions;
\item
  \(G\): relational organization among processes and persons;
\item
  \(\Phi\): one or more meaning maps;
\item
  \(H\): developmental, relational, and learning history;
\item
  \(\pi\): scale, perspective, access map, or aggregation rule;
\item
  \(T\): temporal horizon; and
\item
  \(\mathcal O\): observation model linking latent constructs to report,
  behavior, physiology, context, and third-person assessment.
\end{itemize}

This tuple is a declaration schema, not a theory that every study must
instantiate literally. Its purpose is to expose hidden choices.

The standardized Domain Analytic Frame above is the compact CAP record
for this paper. It also declares the scalar, local--global, allocation,
and translation status required by proposed CU-FP2 and CU-M11--CU-M14.

A claim that ``avoidance is adaptive,'' for example, is incomplete until
the threat regime, available alternatives, protected functions, social
consequences, and horizon are specified. A claim that a person has an
``incoherent identity'' is incomplete until the identity criterion,
cultural expectations, contexts, and practical stakes of the judgment
are stated.

\subsection{2.2 Meaning space}\label{meaning-space}

\textbf{{[}Definition{]} A psychological meaning space} is a
model-specific space whose dimensions represent empirically or
clinically relevant relations among perception, affect, memory,
identity, motivation, embodiment, values, social belonging, and
anticipated futures.

There is no universal meaning space. One study may model appraisals of
threat and controllability. Another may model narrative themes of agency
and communion. A third may use ecological momentary assessment of mood,
sleep, conflict, rumination, and support. A fourth may formalize action
affordances in a sensorimotor task.

The central methodological requirement is that dimensions be justified
by the target phenomenon and observation process. A low-dimensional
visualization can be useful while omitting precisely the relation that
matters to the person.

\subsection{2.3 Meaning maps}\label{meaning-maps}

\textbf{{[}Definition{]} A meaning map} is an organized, revisable
family of relations through which events, persons, sensations, memories,
possibilities, and actions acquire significance relative to
expectations, identity, goals, relationships, values, and futures.

The notation

\[\adjustbox{max width=\linewidth}{$\displaystyle 
\Phi_t:
(\text{event},\text{context},\text{history},\text{possible action})
\longmapsto
(\text{expectation},\text{affordance},\text{significance})
$}\]

is schematic. It does not imply one centralized map, one neural
location, or a fully explicit symbolic representation.

At least five forms should be distinguished:

\begin{itemize}
\tightlist
\item
  \textbf{propositional maps:} beliefs and explicit explanations;
\item
  \textbf{predictive maps:} learned expectations about what follows
  what;
\item
  \textbf{affective maps:} patterns of attraction, aversion, safety, and
  salience;
\item
  \textbf{embodied maps:} action readiness, posture, interoceptive
  expectation, and sensorimotor habit;
\item
  \textbf{narrative and social maps:} roles, identities, scripts, and
  anticipated recognition; and
\item
  \textbf{commitment maps:} values, vows, sacred meanings, and forms of
  life one is prepared to enact.
\end{itemize}

A person can explicitly believe that a room is safe while their body
remains prepared for threat. They can endorse self-compassion while
automatically interpreting error as rejection. They can hold a religious
hope without claiming to possess descriptive certainty. Such divergences
are not anomalies; they are ordinary properties of layered psychological
organization.

\subsection{2.4 The psychological coherence
profile}\label{the-psychological-coherence-profile}

Psychological organization is represented here by a profile:

\[\adjustbox{max width=\linewidth}{$\displaystyle 
\mathbf P_{\Psi}
=
\langle
p_{\mathrm{Ad}},
p_{\mathrm{Af}},
p_{\mathrm{Fr}},
p_{\mathrm{In}},
p_{\mathrm{Rc}},
p_{\mathrm{Ga}},
p_{\mathrm{Br}},
p_{\mathrm{Mp}},
p_{\mathrm{Qr}}
\rangle .
$}\]

The labels denote candidate dimensions:

\begin{itemize}
\tightlist
\item
  \(p_{\mathrm{Ad}}\): \textbf{descriptive adequacy}---calibration to
  evidence and reality;
\item
  \(p_{\mathrm{Af}}\): \textbf{affect tolerance and
  regulation}---capacity to experience, identify, modulate, and use
  affect without automatic collapse or exclusion;
\item
  \(p_{\mathrm{Fr}}\): \textbf{regulatory flexibility}---ability to vary
  strategies with context and feedback;
\item
  \(p_{\mathrm{In}}\): \textbf{identity and narrative
  integration}---continuity without rigid self-protection;
\item
  \(p_{\mathrm{Rc}}\): \textbf{relational capacity}---connection,
  boundaries, reciprocity, trust, repair, and recognition;
\item
  \(p_{\mathrm{Ga}}\): \textbf{agency}---ability to initiate, inhibit,
  choose, revise, and act in relation to goals and values;
\item
  \(p_{\mathrm{Br}}\): \textbf{embodied regulation}---sleep, arousal,
  interoceptive use, sensorimotor organization, and bodily safety;
\item
  \(p_{\mathrm{Mp}}\): \textbf{meaning and purpose}---experienced
  comprehension, direction, significance, and mattering; and
\item
  \(p_{\mathrm{Qr}}\): \textbf{repair capacity}---ability to register
  disruption, seek or receive support, learn, reorganize, and recover or
  transform.
\end{itemize}

The profile is extensible. It is not a diagnostic instrument, and no
dimension is assumed to be measurable by one universal scale.

\subsubsection{2.4.1 Partial order and
dominance}\label{partial-order-and-dominance}

For two profiles \(p\) and \(q\), define weak dominance relative to a
declared set of dimensions \(J\):

\[\adjustbox{max width=\linewidth}{$\displaystyle 
p\succeq_J q
\quad\Longleftrightarrow\quad
p_j\ge q_j\quad\text{for all }j\in J.
$}\]

If each profile is better on at least one relevant dimension, they are
non-dominating. A scalar ranking then requires additional normative,
clinical, or decision assumptions.

\subsubsection{2.4.2 Scalarization is a decision, not a
discovery}\label{scalarization-is-a-decision-not-a-discovery}

A model may define

\[\adjustbox{max width=\linewidth}{$\displaystyle 
S_w(p)=\sum_j w_j p_j
$}\]

for a declared purpose. But the weights \(w_j\), normalizations, floors,
and excluded dimensions carry substantive commitments. If changing
reasonable weights reverses a ranking, the correct conclusion is not
that one weighting discovered the true order. The correct conclusion is
that the profiles encode a tradeoff.

\subsubsection{2.4.3 Protected floors}\label{protected-floors}

Some decisions require constraints rather than compensation. A treatment
plan should not allow a small gain in symptom score to compensate for
suicidality, loss of decision capacity, intolerable adverse effects, or
violation of consent. A workplace should not offset severe
discrimination with average productivity.

A constrained comparison may therefore require

\[\adjustbox{max width=\linewidth}{$\displaystyle 
p_j\ge \ell_j
\qquad
\text{for every protected dimension }j\in J_K.
$}\]

The floors must be declared and justified. They are not supplied by
mathematics alone.

\subsection{2.5 Integration load and regulatory
capacity}\label{integration-load-and-regulatory-capacity}

Psychological strain should not be treated as the inverse of one hidden
integration quantity. It is modeled here as a relation among demands,
uncertainty, conflict, support, bodily condition, history, available
action, power, and horizon.

Let \(d_t\) be a vector of demands: threat, uncertainty, role conflict,
sensory load, grief, cognitive complexity, economic pressure,
caregiving, pain, social evaluation, or moral injury. Let

\[\adjustbox{max width=\linewidth}{$\displaystyle 
\mathcal C_t=\mathcal C(x_t,E_t,H,\pi,T)
$}\]

be the set of demand configurations the person can currently meet
without violating protected functions.

\textbf{{[}Definition{]} Integration overload} occurs when

\[\adjustbox{max width=\linewidth}{$\displaystyle 
d_t\notin\mathcal C_t.
$}\]

This relation is model-relative. It does not imply that distress is a
defect. Distress may be an appropriate signal of danger, loss,
injustice, overload, or conflict. Nor does the absence of distress prove
health; numbing, mania, dissociation, intoxication, or enforced
compliance can reduce reported discomfort while increasing danger.

\subsection{2.6 Local, global, short, and long
horizons}\label{local-global-short-and-long-horizons}

The original local--global principle remains one of the framework's
strongest insights when made precise.

\textbf{{[}Bridge principle{]} A psychological strategy may stabilize
one process, context, relationship, scale, or time interval while
destabilizing another.}

Examples include:

\begin{itemize}
\tightlist
\item
  avoidance that lowers immediate fear while shrinking future action;
\item
  emotional suppression that permits role performance while weakening
  intimacy;
\item
  rigid identity that preserves continuity while blocking correction;
\item
  masking that increases social acceptance while raising exhaustion or
  self-alienation;
\item
  medication that reduces one symptom while producing adverse effects;
\item
  confrontation that increases short-term distress but supports learning
  under adequate safety;
\item
  spiritual practice that widens tolerance or, alternatively, becomes a
  vehicle for bypass;
\item
  group belonging that supports regulation while demanding factual or
  moral conformity.
\end{itemize}

No strategy is classified by its immediate effect alone.

\subsubsection{2.6.1 Reader-facing horizon-reversal
construction}\label{reader-facing-horizon-reversal-construction}

The released technical supplement includes a deterministic four-variable
construction with three analyst-defined regulation regimes. It is a
pedagogical contrast, not a treatment model or clinical finding. The
construction declares acute distress, sensitization, valued functioning,
and regulatory capacity, then compares suppression-dominant,
supported-engagement, and unbuffered-confrontation regimes under one
common probe.

\begin{longtable}[]{@{}
  >{\raggedright\arraybackslash}p{(\columnwidth - 8\tabcolsep) * \real{0.1579}}
  >{\raggedleft\arraybackslash}p{(\columnwidth - 8\tabcolsep) * \real{0.2105}}
  >{\raggedleft\arraybackslash}p{(\columnwidth - 8\tabcolsep) * \real{0.2105}}
  >{\raggedleft\arraybackslash}p{(\columnwidth - 8\tabcolsep) * \real{0.2105}}
  >{\raggedleft\arraybackslash}p{(\columnwidth - 8\tabcolsep) * \real{0.2105}}@{}}
\toprule\noalign{}
\begin{minipage}[b]{\linewidth}\raggedright
Regime
\end{minipage} & \begin{minipage}[b]{\linewidth}\raggedleft
Mean acute distress
\end{minipage} & \begin{minipage}[b]{\linewidth}\raggedleft
Final functioning
\end{minipage} & \begin{minipage}[b]{\linewidth}\raggedleft
Final capacity
\end{minipage} & \begin{minipage}[b]{\linewidth}\raggedleft
Probe peak distress
\end{minipage} \\
\midrule\noalign{}
\endhead
\bottomrule\noalign{}
\endlastfoot
Suppression-dominant & 0.00445 & 0.49283 & 0.59470 & 0.93557 \\
Supported engagement & 0.15415 & 0.88518 & 0.80450 & 0.31532 \\
Unbuffered confrontation & 0.43371 & 0.80464 & 0.56113 & 0.45874 \\
\end{longtable}

Suppression produces the least acute distress but the poorest later
probe response. Supported engagement produces more acute distress while
improving the three later dimensions. The two profiles are therefore
non-dominating. A short-horizon weighting ranks suppression first
(0.93018 versus 0.84015), while a longer-horizon weighting ranks
supported engagement first (0.79090 versus 0.41863). Unbuffered
confrontation exceeds the declared acute-distress floor and is strictly
dominated by supported engagement on the four evaluation coordinates.

The construction establishes only that immediate relief, later
functioning, protected floors, and scalarization can come apart under
declared equations. Its regime definitions deliberately reward supported
learning, so robustness within the construction is not clinical
evidence. Full equations, sensitivity checks, alternative initial
states, and revision boundaries are in the released Psychology technical
supplement.

\subsection{2.7 Observation and
inference}\label{observation-and-inference}

Psychological constructs are not read directly from one channel.
Relevant evidence may include:

\begin{itemize}
\tightlist
\item
  first-person report;
\item
  behavior and task performance;
\item
  ecological momentary assessment;
\item
  longitudinal course;
\item
  physiological and neural measures;
\item
  collateral report;
\item
  intervention and ablation;
\item
  cultural and material context;
\item
  and the person's own goals and interpretation.
\end{itemize}

Disagreement among channels is data. It can reflect limited insight,
measurement error, context variation, stigma, demand characteristics,
dissociation, observer bias, or genuinely multidimensional organization.

The observation model \(\mathcal O\) must therefore be declared. A
self-report questionnaire is not the construct itself. A clinician
rating is not neutral access to the person. A physiological marker is
not a direct readout of meaning.

\subsection{2.8 Translation cards}\label{translation-cards}

Every cross-domain use of a CU term should answer eight questions:

\begin{enumerate}
\def\labelenumi{\arabic{enumi}.}
\tightlist
\item
  What is the source construct?
\item
  What is the psychological target construct?
\item
  Which relation is preserved?
\item
  Which new constraints enter?
\item
  What information is lost?
\item
  What evidence supports the translation?
\item
  Which alternatives remain viable?
\item
  What would make the translation fail?
\end{enumerate}

For example, translating biological viability into psychological agency
preserves the relevance of protected functions and trajectories, but it
adds representation, social meaning, conscious experience, and
person-specific goals. It loses any warrant to infer that biological
persistence and psychological flourishing are the same objective.

\begin{center}\rule{0.5\linewidth}{0.5pt}\end{center}

\section{3. Perception, memory, and the construction of
meaning}\label{perception-memory-and-the-construction-of-meaning}

\subsection{3.1 Perception is active without being
arbitrary}\label{perception-is-active-without-being-arbitrary}

Perception is not a passive copy of the world. Organisms and persons
sample, organize, and interpret sensory input through prior learning,
attention, bodily state, task, and expectation. Predictive-processing
and active-inference approaches provide influential model families for
this interaction, but this paper does not treat one variational
formulation as the universal mathematics of mind (Clark, 2013; Friston,
2010; Hohwy, 2013).

The fact that perception is constructed does not imply that every
construction is equally adequate. Perceptual systems remain answerable
to environmental structure through error, action, intersubjective
checking, and consequence. Hallucinations, illusions, culturally learned
categories, and ordinary perception differ in mechanism and evidential
relation even though all involve active organization.

\textbf{{[}Bridge principle{]} Perceptual coherence concerns the
coordination of sensory evidence, prior expectation, attention, bodily
state, and action---not agreement with expectation at any cost.}

A system that suppresses all disconfirming input may be internally
stable and descriptively poor. Descriptive adequacy therefore remains an
independent profile dimension.

\subsection{3.2 Attention as selection under limited
capacity}\label{attention-as-selection-under-limited-capacity}

Attention changes which elements of a situation can influence memory,
appraisal, and action. Threat can narrow attention adaptively; sustained
narrowing can also maintain anxiety, prejudice, or trauma-related
monitoring. Curiosity can widen sampling; overload can make widening
impossible.

Attentional selection should be evaluated relative to task and cost. A
narrow focus may be beneficial in expert performance or monotropic
interest and disabling in a rapidly changing social environment. A broad
focus may aid exploration while impairing completion. The psychological
question is not which style is globally ``more coherent,'' but whether
the selection regime serves the person's goals and protected functions
within the relevant context.

\subsection{3.3 Memory as constructive
continuity}\label{memory-as-constructive-continuity}

Memory is not a neutral archive. Retrieval is reconstructive,
cue-dependent, and vulnerable to misinformation and current goals
(Bartlett, 1932; Loftus, 2005; Schacter, 1996). This does not make
memory unreal. It means that continuity is produced through repeated
reactivation, stabilization, rehearsal, social narration, bodily cueing,
and revision.

This paper distinguishes:

\begin{itemize}
\tightlist
\item
  \textbf{event memory:} representation of particular episodes;
\item
  \textbf{semantic memory:} generalized knowledge;
\item
  \textbf{procedural memory:} learned skill and habit;
\item
  \textbf{affective memory:} learned salience, safety, and threat
  relations;
\item
  \textbf{autobiographical memory:} remembered material organized around
  a life and self;
\item
  \textbf{social memory:} records distributed through relationships and
  institutions; and
\item
  \textbf{commitment memory:} continuity of values, promises, and
  practices.
\end{itemize}

These forms bind differently to identity. Loss of episodic recall need
not erase character, relationship, or practical identity. Conversely,
preserved facts do not guarantee preserved agency or selfhood.

\subsection{3.4 Memory updating and
reconsolidation}\label{memory-updating-and-reconsolidation}

Retrieved memories can sometimes become modifiable before
restabilization. Experimental reconsolidation research has established
important boundary conditions, especially in animal learning and
selected human paradigms. Its translation into psychotherapy remains an
active, uneven research program rather than a universal mechanism of
change (Elsey, Van Ast, \& Kindt, 2018; Lee, Nader, \& Schiller, 2017).

The publication-safe claim is:

\begin{quote}
\textbf{{[}Conditional result{]} Some therapeutic procedures may alter
the emotional, predictive, or behavioral consequences of a memory when
retrieval occurs under conditions that permit new learning. This does
not imply erasure, guarantee destabilization, or establish one mechanism
across therapies.}
\end{quote}

Nonlinear change remains possible, but its mechanism must be shown case
by case; memory updating does not by itself establish one universal
``phase-leap'' law.

\subsection{3.5 Meaning as a relation, not a
substance}\label{meaning-as-a-relation-not-a-substance}

Meaning is not an amount stored inside an event. It is a relation among
event, person, history, possible action, social world, and future.

An event can be physically similar while psychologically different
because it occupies a different relation in the person's map. A racing
heart during exercise, panic, attraction, illness, or public performance
has different anticipated causes and consequences. The bodily signal
participates in meaning but does not determine it alone.

Research on meaning in life commonly distinguishes comprehension or
coherence, purpose, and significance or mattering (Steger et al., 2006;
King \& Hicks, 2021; Martela \& Steger, 2016). This paper retains those
constructs while preventing terminological confusion: \emph{coherence}
in meaning-in-life research is a person's felt comprehensibility of
life, not the framework's entire cross-domain concept.

\subsection{3.6 Description and commitment inside meaning
maps}\label{description-and-commitment-inside-meaning-maps}

A meaning map can contain both truth-apt descriptions and orienting
commitments.

Consider:

\begin{quote}
``My husband has died.''
\end{quote}

This is a descriptive claim. It answers to evidence.

\begin{quote}
``My life can still be lived in fidelity to our love.''
\end{quote}

This may function principally as a commitment. It is not insulated from
facts or consequences, but its role is not exhausted by predicting an
independently specifiable event.

Confusion arises in both directions. A commitment can be falsely
presented as scientific knowledge. A symbolic or sacred utterance can
also be criticized as though its only function were literal prediction.
The two-register architecture does not separate these domains
absolutely; it makes their traffic explicit.

\subsection{3.7 Narrative identity}\label{narrative-identity}

Narrative identity organizes remembered past, interpreted present, and
imagined future. Themes such as agency, communion, redemption,
contamination, belonging, continuity, and rupture can shape well-being
and action (McAdams, 2001; McAdams \& McLean, 2013). Longitudinal
evidence suggests that narrative themes can explain variance in
depression and well-being beyond broad personality traits, while not
functioning as deterministic causes (Lind et al., 2024).

A coherent narrative is not necessarily accurate, healthy, or complete.
Propaganda, grandiosity, conspiracy, and family myths can offer powerful
narrative closure. Conversely, uncertainty or revision can indicate
honesty and growth.

\textbf{{[}Boundary principle{]} Narrative integration must be assessed
alongside descriptive adequacy, affect tolerance, agency, relationship,
and openness to revision.}

\subsection{3.8 Meaning-map mismatch}\label{meaning-map-mismatch}

A mismatch can occur at several levels:

\begin{itemize}
\tightlist
\item
  a descriptive belief conflicts with evidence;
\item
  an implicit prediction conflicts with explicit belief;
\item
  an embodied response conflicts with current safety;
\item
  a social role conflicts with personal values;
\item
  a commitment conflicts with its consequences;
\item
  a narrative excludes important experience;
\item
  a person's map fits one environment and not another.
\end{itemize}

Not every mismatch is pathology. Some are ordinary features of learning.
Some are signs of plural roles. Some reflect unresolved injustice. Some
become clinically important when they produce sustained distress,
danger, disability, loss of agency, or inability to pursue the person's
aims.

\subsection{3.9 The epistemic role of
emotion}\label{the-epistemic-role-of-emotion}

Emotion can direct attention toward needs, threats, losses, violations,
attachments, and opportunities. It can also mislead when learned
associations, current physiology, strategic manipulation, or prior
trauma make a signal poorly calibrated to the present.

The correct alternative to ``emotion is noise'' is not ``emotion is
infallible data.'' Emotional experience is evidence about the person's
current organization and history. It is not automatically evidence that
the external interpretation attached to it is true.

\subsection{3.10 Memory, testimony, and
humility}\label{memory-testimony-and-humility}

Because memory and interpretation are reconstructive, psychological
inquiry requires humility. But humility does not mean reflexive
disbelief toward first-person testimony. Reports of pain, fear,
violation, identity, or meaning are indispensable evidence about lived
experience, even when causal interpretation remains uncertain.

A responsible framework therefore avoids two errors:

\begin{itemize}
\tightlist
\item
  treating first-person report as transparent and incorrigible access to
  mechanism; and
\item
  treating third-person measurement as the only legitimate form of
  knowledge.
\end{itemize}

The observation model must represent both.

\begin{center}\rule{0.5\linewidth}{0.5pt}\end{center}

\section{4. Affect, motivation, embodiment, and
agency}\label{affect-motivation-embodiment-and-agency}

\subsection{4.1 Affect is
multidimensional}\label{affect-is-multidimensional}

\textbf{Functional stake is not phenomenal valence:} a condition can
matter to regulation, survival, attachment, or an endorsed aim without
thereby being felt as pleasant or unpleasant. Human affect is organized
through interacting dimensions that may include valence, arousal, action
tendency, appraisal, bodily state, social meaning, and temporal
expectation.

Emotion theories disagree about the relative roles of basic affective
systems, appraisal, construction, prediction, and social learning. This
paper does not settle that debate. It adopts a minimal common ground:

\begin{quote}
\textbf{{[}Empirical generalization{]} Emotions coordinate perception,
bodily preparation, attention, learning, communication, and action in
context-sensitive ways.}
\end{quote}

That coordination can be adaptive, biased, conflicted, or dysregulated.

\subsection{4.2 Emotion regulation is not emotion
elimination}\label{emotion-regulation-is-not-emotion-elimination}

Emotion regulation includes processes by which people influence which
emotions occur, when they occur, how they are experienced, and how they
are expressed (Gross, 2002, 2015). Strategies such as reappraisal,
suppression, rumination, acceptance, avoidance, problem solving,
distraction, support-seeking, and exposure have effects that depend on
context, timing, intensity, culture, and goal.

A large 2026 meta-analysis found broad self-reported emotion-regulation
difficulties across many mental disorders while also finding
disorder-specific profiles; rumination, avoidance, suppression, reduced
acceptance, and reduced reappraisal were not identical in distribution
or effect (Clamor, Lincoln, \& Schulze, 2026). This supports a
transdiagnostic role for regulation while warning against one deficit
score.

\textbf{{[}Bridge principle{]} Regulation should be assessed by
flexibility, context fit, cost, and horizon---not by the frequency of
one strategy alone.}

\subsection{4.3 Regulatory flexibility}\label{regulatory-flexibility}

Psychological flexibility is the capacity to remain in contact with
present conditions, vary behavior with context, and act in relation to
values despite difficult private events. Regulatory-flexibility research
likewise emphasizes matching strategies to changing contexts rather than
treating any single strategy as universally healthy (Bonanno \& Burton,
2013). These constructs have substantial transdiagnostic and
treatment-process support, but their measurement and causal status
remain active research questions (Kashdan \& Rottenberg, 2010; Macri \&
Rogge, 2024).

This paper uses \textbf{regulatory flexibility} more broadly than any
one therapeutic school. It includes:

\begin{itemize}
\tightlist
\item
  detecting when a strategy is not working;
\item
  switching among strategies;
\item
  tolerating short-term discomfort when needed;
\item
  accepting limits that cannot be controlled;
\item
  seeking support;
\item
  revising goals;
\item
  and preserving protected functions during change.
\end{itemize}

Flexibility does not mean endless openness. Boundaries, commitments,
routines, and refusal can be flexible responses when they fit the
situation.

\subsection{4.4 Motivation is plural}\label{motivation-is-plural}

People do not pursue one hidden objective. Motivation can be organized
by reward, habit, attachment, threat avoidance, curiosity, competence,
autonomy, belonging, status, care, identity, justice, meaning, and
commitment. Self-determination theory, reinforcement learning,
expectancy-value models, attachment, psychodynamic theory, and
existential psychology each describe different parts of this landscape
(Deci \& Ryan, 2000).

Human motivation is plural. The defensible bounded claim is:

\begin{quote}
\textbf{{[}Empirical generalization{]} People often seek to reduce
threatening discrepancy and restore a workable relation among
expectation, identity, affect, belonging, and action.}
\end{quote}

This tendency is not inherently truth-seeking or health-promoting.
Conspiracy, addiction, compulsive ritual, domination, self-deception,
and rigid ideology can all reduce some discrepancies while creating
others.

Plurality does not make motivation directionless. A person registers not
only where they stand on variables that matter to them but how those
variables are changing and what candidate actions are expected to do.
\emph{Consciousness} §8.4 develops this as \textbf{directional
evaluation} (CU-V8): a conditional phasic hypothesis about change in
endogenous better--worse organization together with a tonic hypothesis
about standing relative to protected margins. Where valence is present,
improving direction may be felt as relief or satisfaction, worsening as
distress, and disagreement among goals or horizons as ambivalence. This
is not one universal gradient and not every motivation is truth-seeking
or health-promoting; short-horizon relief can accompany long-horizon
narrowing.

\subsection{4.5 Needs, values, and
commitments}\label{needs-values-and-commitments}

Needs are not interchangeable with preferences. Sleep, safety,
nutrition, attachment, autonomy, competence, recognition, and bodily
integrity can impose constraints on functioning. Values and commitments
orient action across time. Goals translate them into projects.

A person may value intimacy while avoiding vulnerability, value health
while using a substance for immediate relief, or value truth while
protecting a cherished identity. Such conflict does not imply hypocrisy
by itself. It creates a problem of regulation, learning, and sometimes
moral responsibility.

\subsection{4.6 Embodied regulation}\label{embodied-regulation}

Psychological processes are carried through bodies. Sleep, pain,
illness, endocrine state, autonomic arousal, interoception, movement,
medication, nutrition, disability, and environmental affordances change
what can be perceived, tolerated, remembered, and enacted.

Interoception is not a single accuracy variable. Self-report, objective
detection, attention, interpretation, and metacognitive awareness can
dissociate. A 2024 meta-analysis found that anxiety was associated with
several negative self-reported relations to bodily signals, while not
supporting a simple equation between anxiety and objective interoceptive
accuracy (Clemente, Murphy, \& Murphy, 2024). Reviews of bodily
self-awareness similarly emphasize distributed brain--body networks
rather than one interoceptive center (Candia-Rivera et al., 2024).

The technical claim is therefore modest:

\begin{quote}
\textbf{{[}Bridge principle{]} Embodied state constrains psychological
regulation, and psychological learning constrains how bodily signals are
attended to and interpreted.}
\end{quote}

\subsection{4.7 Insight is not
embodiment}\label{insight-is-not-embodiment}

A person can understand a pattern conceptually while remaining unable to
enact a different response under arousal. This is not evidence that
insight is useless. It shows that propositional revision does not
automatically update procedural, affective, relational, and bodily maps.

Therapeutic work may therefore require repeated action, new relational
experience, sleep and recovery, environmental change, medication,
movement, exposure, rehearsal, or practical support. No single route is
universal.

\subsection{4.8 Agency}\label{agency}

\textbf{{[}Definition{]} Psychological agency} is the graded capacity to
initiate, inhibit, sustain, revise, and coordinate action in relation to
perceived options, goals, values, commitments, and consequences.

Agency depends on more than intention. It can be constrained by fatigue,
threat, compulsion, addiction, pain, poverty, coercion, cognitive
impairment, social exclusion, and unavailable alternatives.

A person can retain moral worth while agency is diminished. They can
also exercise causal power without reflective authorship. The framework
therefore distinguishes:

\begin{itemize}
\tightlist
\item
  \textbf{behavioral control};
\item
  \textbf{choice among represented options};
\item
  \textbf{reflective endorsement};
\item
  \textbf{diachronic authorship};
\item
  \textbf{responsibility under fair conditions}; and
\item
  \textbf{political or institutional power}.
\end{itemize}

\subsection{4.9 Addiction and
compulsion}\label{addiction-and-compulsion}

Addiction and compulsion illustrate local--global tension without
reducing either condition to it. Repeated behavior can provide immediate
relief, reward, identity, belonging, or withdrawal avoidance while
altering attention, learning, health, relationship, and future choice.

The relevant profile may include:

\begin{itemize}
\tightlist
\item
  cue sensitivity;
\item
  craving and withdrawal;
\item
  habit strength;
\item
  alternative rewards;
\item
  executive control;
\item
  stress exposure;
\item
  social network;
\item
  shame;
\item
  meaning;
\item
  and material access to treatment.
\end{itemize}

Calling the pattern a ``coherence trap'' is permissible only as an
interpretation after the mechanisms are specified.

\subsection{4.10 Affect and morality}\label{affect-and-morality}

Affective reactions contribute to moral attention and motivation, but
they do not settle justification. Disgust, loyalty, empathy, anger,
fear, and compassion can disclose morally relevant features and can be
biased by group boundaries, misinformation, or self-interest.

Psychology can study how moral judgments are formed and distorted. It
cannot derive the complete content of morality from whichever responses
are most stable or strongly felt.

\begin{center}\rule{0.5\linewidth}{0.5pt}\end{center}

\section{5. Identity, personality, and
self-deception}\label{identity-personality-and-self-deception}

\subsection{5.1 Identity is plural and
decision-relative}\label{identity-is-plural-and-decision-relative}

A person can be identified through body, biography, memory,
relationship, legal standing, values, roles, commitments, and patterns
of action. These criteria overlap without being identical.

Following Foundations v0.10, an identity claim must state:

\begin{itemize}
\tightlist
\item
  the entity being tracked;
\item
  the identity-relevant constraints;
\item
  the admissible transformations;
\item
  the temporal horizon;
\item
  and the practical decision the identity judgment serves.
\end{itemize}

The question ``Am I still the same person?'' can therefore concern legal
continuity, autobiographical memory, moral character, relational trust,
spiritual orientation, or felt selfhood. There need not be one answer
across all uses.

\subsection{5.2 Self-model and
self-maintenance}\label{self-model-and-self-maintenance}

A self-model organizes information about body, capacities, roles, social
standing, values, history, and possible action. It is not identical with
the whole person. Nor does self-modeling automatically imply egoic
self-maintenance.

\emph{Consciousness} v0.10 defines ego as ``a persistent
self-maintaining organization that treats the continuation, status,
integrity, or narrative of a modeled self as an object of defense.''
Healthy self-protection, boundaries, and continuity do not thereby
become pathological ego.

Psychological maturity cannot mean eliminating every self-model. It may
involve using self-representation without making every correction, loss,
or disagreement into a threat to existence.

\subsection{5.3 Personality as patterned
organization}\label{personality-as-patterned-organization}

Personality traits describe relatively stable differences in affect,
behavior, cognition, and motivation. Traits are not rigid essences; they
vary across situations and can change across the lifespan. Narrative
identity, characteristic adaptations, roles, and culture add levels not
captured by broad traits alone.

This paper treats personality as a patterned profile of regulatory
tendencies rather than a single ``coherence strategy.'' A trait can be
advantageous in one environment and costly in another. Conscientiousness
can support reliability and become rigidity. Openness can support
exploration and undermine completion. Agreeableness can support care and
weaken boundaries. Neuroticism can increase distress and heighten threat
detection.

\subsection{5.4 Self-concept clarity and
revision}\label{self-concept-clarity-and-revision}

Self-concept clarity refers to the perceived definition, consistency,
and stability of beliefs about oneself (Campbell et al., 1996). Lower
clarity is associated with several forms of distress, but high clarity
is not always truth or health. A rigid, inaccurate self-concept can be
very clear.

The relevant distinction is between:

\begin{itemize}
\tightlist
\item
  \textbf{clarity:} how definite the self-model feels;
\item
  \textbf{adequacy:} how well it fits evidence and lived experience;
\item
  \textbf{flexibility:} whether it can revise;
\item
  \textbf{integration:} whether conflicting roles and histories can be
  held;
\item
  \textbf{authorship:} whether commitments are reflectively owned; and
\item
  \textbf{social recognition:} whether others permit the identity to be
  lived.
\end{itemize}

\subsection{5.5 Narrative continuity and
rupture}\label{narrative-continuity-and-rupture}

Life stories can preserve continuity through change. They can also
imprison a person in a fixed role: victim, savior, failure, chosen one,
outsider, strong one, good child, rational one, or damaged person.

A rupture may initially reduce coherence while enabling a more adequate
identity. Divorce, conversion, deconversion, disability, migration,
grief, parenthood, recovery, and moral injury can require a period in
which no stable story yet fits.

Repair is therefore not always restoration. Sometimes it is the
construction of a continuity capable of acknowledging that something
irreversible occurred.

\subsection{5.6 Defining self-deception}\label{defining-self-deception}

\textbf{{[}Definition{]} Self-deception is a motivated pattern in which
attention, memory, interpretation, affect, or standards are selectively
organized to protect a threatened commitment, identity, attachment,
interest, or policy from reasons that would otherwise require revision.}

This definition has five components:

\begin{enumerate}
\def\labelenumi{\arabic{enumi}.}
\tightlist
\item
  \textbf{relevant reasons:} evidence or experience bears on the
  protected conclusion;
\item
  \textbf{asymmetry:} favorable and unfavorable information are
  processed by different standards;
\item
  \textbf{protective function:} the asymmetry stabilizes an identity,
  attachment, benefit, or policy;
\item
  \textbf{counterfactual sensitivity:} the pattern would change if the
  same evidence supported the preferred conclusion, or if the identity
  stake were removed; and
\item
  \textbf{persistence against correction:} the process resists
  accessible opportunities for revision beyond ordinary uncertainty or
  limited attention.
\end{enumerate}

No one component is sufficient in every case. The profile guides inquiry
rather than functioning as an accusation algorithm.

\subsubsection{5.6.1 Correctability and the non-culpable
case}\label{correctability-and-the-non-culpable-case}

Self-deception is the motivated form of a wider loss of correctability:
evidence, meaning, or standards are asymmetrically reorganized to
protect a stake. Trauma and protective compartmentalization may also
narrow correctability, but often non-culpably. A model that once tracked
real danger can outlive its conditions without being dishonest. Healing
therefore restores reality contact, tolerance, and revision capacity
without turning survival adaptations, ordinary error, uncertainty, or
disagreement into blame.

\subsection{5.7 What self-deception is
not}\label{what-self-deception-is-not}

Self-deception must be distinguished from ordinary inconsistency and
from cognitive dissonance as a descriptive state (Festinger, 1957). It
must also be distinguished from:

\begin{itemize}
\tightlist
\item
  \textbf{error:} a false belief held without motivated asymmetry;
\item
  \textbf{uncertainty:} insufficient evidence or unresolved
  interpretation;
\item
  \textbf{confabulation:} production of an explanation without conscious
  access to the actual cause;
\item
  \textbf{repression or dissociation:} limited access to material, which
  may or may not be motivated in the relevant sense;
\item
  \textbf{adaptive deferral:} postponing processing while conditions are
  unsafe;
\item
  \textbf{privacy:} declining to disclose;
\item
  \textbf{value disagreement:} endorsing a different commitment; and
\item
  \textbf{ordinary inconsistency:} failing to integrate every belief and
  action.
\end{itemize}

The moral importance of self-deception does not license casual
mind-reading.

\subsection{5.8 Evidence for
self-deception}\label{evidence-for-self-deception}

Evidence strengthens when several patterns converge:

\begin{itemize}
\tightlist
\item
  selective search or exposure;
\item
  asymmetric skepticism;
\item
  memory or attention biased toward protection;
\item
  standard shifting across equivalent cases;
\item
  avoidance of disconfirming situations;
\item
  affective escalation when identity is threatened;
\item
  post hoc reconstruction;
\item
  strategic ambiguity;
\item
  and institutional reward for not seeing.
\end{itemize}

The same behavior can have multiple explanations. A person who avoids a
conversation may be self-deceiving, overwhelmed, unsafe, exhausted, or
reasonably protecting a boundary. Context and alternatives matter.

\subsection{5.9 Self-deception and the two
registers}\label{self-deception-and-the-two-registers}

Commitments do not become immune merely because they are not ordinary
empirical hypotheses. A religious, political, therapeutic, or relational
commitment becomes unwarranted when it requires factual falsehood,
selective exclusion of suffering, destruction of agency, or systematic
evasion of consequence.

Conversely, acknowledging that a commitment is not known as fact can
reduce self-deception. One can live through trust without claiming
certainty.

\subsection{5.10 Shame, guilt, and repair}\label{shame-guilt-and-repair}

Shame organizes the whole self as defective; guilt more often concerns
an action or responsibility. The distinction is neither absolute nor
universally beneficial. Shame can motivate concealment and attack; guilt
can motivate repair; either can become excessive, manipulated, or
culturally distorted.

A repair-oriented response asks:

\begin{itemize}
\tightlist
\item
  What happened?
\item
  Who was affected?
\item
  What responsibility is mine?
\item
  What conditions shaped the action?
\item
  What can be repaired?
\item
  What boundary or accountability is required?
\item
  What must change so repetition is less likely?
\end{itemize}

Psychological integration is not achieved by dissolving responsibility
into self-compassion or by turning responsibility into permanent
self-condemnation.

\subsection{5.11 Non-egoic development}\label{non-egoic-development}

The spiritual and psychological aspiration toward non-egoic life should
not be confused with weak identity, passivity, compliance, or lack of
boundaries. A person may become less organized around defending status
and narrative while becoming more capable of action, truth, care, and
refusal.

The empirical question concerns capacities: reduced defensive
reactivity, increased decentering, greater flexibility, less compulsive
self-reference, stronger compassion, and preserved agency. The
metaphysical interpretation belongs to Section 13.

\begin{center}\rule{0.5\linewidth}{0.5pt}\end{center}

\section{6. Attachment, sexuality, intimacy, and
co-regulation}\label{attachment-sexuality-intimacy-and-co-regulation}

\subsection{6.1 Attachment as a regulatory
relationship}\label{attachment-as-a-regulatory-relationship}

Attachment theory begins from the observation that proximity to
particular others regulates threat, exploration, and security,
especially in early development (Bowlby, 1969; Ainsworth et al., 1978).
Later work extends co-regulation into adult partnership, separation,
grief, and recovery (Sbarra \& Hazan, 2008).

This paper treats attachment as an important bridge between biology,
consciousness, psychology, and social life. It rejects the reduction of
attachment to one fixed style or one childhood cause.

\textbf{{[}Empirical bridge{]} Attachment relationships can provide
external regulation through which affect tolerance, prediction, trust,
exploration, self-regulation, and expectations of care develop.}

These capacities are influenced by temperament, culture,
neurodivergence, adversity, later relationships, and material
conditions. Attachment classifications summarize patterns under
particular methods; they do not exhaust a person.

\subsection{6.2 Social baseline and distributed
effort}\label{social-baseline-and-distributed-effort}

Humans often regulate risk and effort through expected access to others.
Social baseline theory proposes that familiar social resources alter the
perceived cost of action and threat (Coan \& Sbarra, 2015). The broader
claim is compatible with evidence that relationships affect health,
stress, and behavior, while mechanisms and causal direction vary.

A person's apparent ``self-regulation'' may therefore be jointly
produced. Sleep, childcare, income, transportation, emotional support,
shared memory, conflict repair, and practical help distribute regulatory
labor.

This makes two errors possible:

\begin{itemize}
\tightlist
\item
  judging dependence as immaturity when interdependence is normal; and
\item
  romanticizing dependence when the relation is coercive, exploitative,
  or disabling.
\end{itemize}

\subsection{6.3 Attachment strategies as conditional
adaptations}\label{attachment-strategies-as-conditional-adaptations}

Anxious, avoidant, disorganized, or secure patterns should not be
treated as moral types. They may represent learned expectations about
availability, danger, closeness, and autonomy. A strategy that was
adaptive in one relational environment can become costly in another.

For example:

\begin{itemize}
\tightlist
\item
  heightened monitoring can detect inconsistent care and produce
  exhaustion;
\item
  emotional distance can preserve functioning under intrusion and block
  later intimacy;
\item
  compliance can protect a child and expose an adult to exploitation;
\item
  self-reliance can preserve agency and prevent receiving help;
\item
  rapid trust can enable connection and increase vulnerability.
\end{itemize}

The question is how the strategy interacts with present conditions and
the person's aims.

\subsection{6.4 Intimacy as negotiated mutual
regulation}\label{intimacy-as-negotiated-mutual-regulation}

Intimacy joins disclosure, recognition, trust, bodily presence,
autonomy, dependence, and risk. It is not maximal merging. A
relationship can be close while preserving difference and boundaries.

\textbf{{[}Bridge principle{]} Relational coherence requires enough
shared understanding and mutual responsiveness for coordination,
together with enough differentiation for each person to retain agency,
privacy, and revision.}

Too little coordination produces chronic uncertainty and misattunement.
Too little differentiation produces control, enmeshment, or erasure.

\subsection{6.5 Sexuality}\label{sexuality}

Sexuality can involve pleasure, identity, attachment, reproduction,
exploration, power, vulnerability, cultural meaning, and bodily
regulation. No one function defines it. No one integration principle
defines sexuality because sexuality can integrate, destabilize, express
care, enact harm, provide relief, create conflict, or remain neutral
depending on context.

A psychologically responsible analysis keeps separate:

\begin{itemize}
\tightlist
\item
  desire;
\item
  consent;
\item
  capacity;
\item
  attachment;
\item
  identity;
\item
  behavior;
\item
  fantasy;
\item
  reproduction;
\item
  social recognition;
\item
  and moral or religious commitment.
\end{itemize}

\subsection{6.6 Consent and agency}\label{consent-and-agency}

Consent is not inferred from the absence of resistance or from apparent
relational harmony. It requires sufficient capacity, information,
voluntariness, revocability, and freedom from coercion. Power
asymmetries, dependency, intoxication, fear, disability, age, and
institutional control can alter what consent means in practice.

The framework's language of integration cannot override a person's
refusal. A relationship is not made coherent by forcing one party to
absorb the other's needs.

\subsection{6.7 Conflict and repair}\label{conflict-and-repair}

Relational conflict is not automatically failure. It can disclose
incompatible expectations, injuries, power imbalances, or needed
differentiation. Repair may involve acknowledgment, apology,
restitution, changed behavior, boundary, separation, or acceptance that
trust will not be restored.

A relationship's repair capacity is not measured by how quickly visible
harmony returns. Premature reconciliation can protect the appearance of
coherence while preserving danger.

\subsection{6.8 Grief and continuing
bonds}\label{grief-and-continuing-bonds}

Grief can disrupt attention, sleep, identity, bodily regulation,
narrative, motivation, and social role. It is not one trajectory and not
inherently pathological. Continuing bonds with the dead can coexist with
adaptation; forgetting is not required for recovery.

Meaning-making can support grief, but hope must not be used to cancel
loss. The two widows in \emph{Denying Knowledge to Make Room for
Commitment} (\emph{Denying Knowledge}, Rader 2026) illustrate the
distinction: identical doctrinal language can organize either avoidance
or integration depending on how it relates to grief, relationship,
truth, and future action.

\begin{center}\rule{0.5\linewidth}{0.5pt}\end{center}

\section{7. Stress, trauma, dissociation, and
resilience}\label{stress-trauma-dissociation-and-resilience}

\subsection{7.1 Terms that must be
separated}\label{terms-that-must-be-separated}

The word \emph{trauma} is used for an event, an injury, a memory, a
nervous-system state, a diagnosis, a cultural history, and a metaphor
for disruption. The analysis requires separate terms.

\begin{itemize}
\tightlist
\item
  A \textbf{potentially traumatic event} is an exposure with serious
  threat, violation, injury, death, coercion, or overwhelming
  helplessness.
\item
  A \textbf{traumatic response} is a pattern of psychological and bodily
  effects following such exposure.
\item
  \textbf{Post-traumatic stress disorder} is a diagnostic syndrome with
  specified criteria.
\item
  \textbf{Complex PTSD} is a distinct ICD-11 category combining core
  PTSD symptoms with disturbances in self-organization.
\item
  \textbf{Dissociation} includes disruptions of integration among
  consciousness, memory, identity, perception, or embodiment and is not
  synonymous with trauma.
\item
  \textbf{Moral injury} concerns enduring psychological, social, or
  spiritual consequences of perpetrating, failing to prevent,
  witnessing, or being betrayed around events that violate deeply held
  moral expectations; it is not identical with PTSD.
\end{itemize}

The 2024 ICD-11 clinical manual includes complex PTSD, prolonged grief
disorder, lifespan guidance, cultural guidance, and dimensional
elements. Diagnostic systems remain tools, not complete causal theories
(World Health Organization, 2024).

\subsection{7.2 Heterogeneous
trajectories}\label{heterogeneous-trajectories}

Exposure does not determine one outcome. Reviews of longitudinal
trajectory studies commonly identify resilient, recovery, chronic, and
delayed patterns, with substantial variation across populations,
measurements, and methods (Galatzer-Levy, Huang, \& Bonanno, 2018). New
complex-systems models formalize resilience as the stability and
recoverability of symptom networks while remaining proof-of-principle
rather than established clinical diagnostics (Lunansky et al., 2025).

Trauma is not one universal ``coherence catastrophe.''

A better statement is:

\begin{quote}
\textbf{{[}Empirical generalization{]} Potentially traumatic events can
disrupt or reorganize affect regulation, bodily prediction, attention,
memory, identity, relationship, agency, meaning, and trust, but the
pattern and duration of change vary substantially.}
\end{quote}

\subsection{7.3 Protective responses}\label{protective-responses}

Fight, flight, freezing, submission, numbing, narrowed attention,
avoidance, dissociation, and compulsive control can reduce danger or
overload under some conditions. Their later cost does not prove that
they were irrational when formed.

This paper calls this \textbf{protective compartmentalization} when
access or coordination among processes is reduced to keep overwhelming
material from dominating the whole system.

Protective compartmentalization can:

\begin{itemize}
\tightlist
\item
  permit survival and role performance;
\item
  limit conscious access to unbearable material;
\item
  reduce immediate arousal;
\item
  preserve attachment under dependency;
\item
  and delay processing until safety exists.
\end{itemize}

It can also:

\begin{itemize}
\tightlist
\item
  impair memory continuity;
\item
  restrict emotional range;
\item
  increase somatic distress;
\item
  fragment identity;
\item
  impair intimacy;
\item
  and make present safety feel like past danger.
\end{itemize}

Its status is therefore context- and horizon-relative.

\subsection{7.4 Trauma and memory}\label{trauma-and-memory}

Traumatic memories can involve vivid intrusion, sensory fragments,
avoidance, overgeneralization, shame, and altered expectations. No
single memory mechanism explains every case. Ordinary autobiographical
memory, fear learning, contextual processing, sleep, social narration,
and repeated retrieval can all contribute.

Claims that ``the body stores trauma'' should be translated carefully.
Bodily responses and procedural expectations can persist without a
propositional narrative, but this does not mean that every tissue
contains a literal episodic record.

\subsection{7.5 Trauma and identity}\label{trauma-and-identity}

An event can alter what the world, self, and future appear to permit. It
can produce contamination narratives, loss of trust, shame,
helplessness, survivor identity, moral rupture, or a sense that one's
former life no longer belongs to oneself.

Identity reconstruction is not automatically positive growth. Some
people find new commitments; some recover without a redemptive story;
some remain injured; some are pressured to produce meaning for others.
No suffering is justified by the possibility that someone may later
grow.

\subsection{7.6 Complex and developmental
trauma}\label{complex-and-developmental-trauma}

Repeated interpersonal threat during dependence can affect attachment,
emotion regulation, self-concept, dissociation, bodily safety, and
relationships. Yet ``developmental trauma'' is not a license to infer a
hidden history from broad symptoms, nor to treat every difficulty as the
result of childhood harm.

Current reviews support several established psychotherapies for PTSD and
complex presentations while emphasizing heterogeneity, modest evidence
certainty for some comparisons, and unresolved questions about
phase-based versus direct trauma-focused treatment (Billings \&
Nicholls, 2025; Schaug et al., 2025).

\subsection{7.7 Dissociation}\label{dissociation}

Dissociation may involve depersonalization, derealization, amnesia,
identity discontinuity, absorption, or altered access among processes.
It can be transient or persistent, adaptive or disabling, trauma-related
or associated with other conditions.

A generic integration score cannot distinguish these forms. Assessment
must specify phenomenology, triggers, function, impairment, and
alternatives such as neurological conditions, sleep disruption,
substance effects, psychosis, imaginative absorption, or culturally
normative experience.

\subsection{7.8 Resilience and repair}\label{resilience-and-repair}

\textbf{{[}Definition{]} Psychological repair capacity} is the ability
of a person-in-context to register disruption, access resources, revise
expectations, restore or transform agency, reconnect with body and
relationship, and widen the range of experience that can be integrated
without violating protected functions.

Repair is not:

\begin{itemize}
\tightlist
\item
  invulnerability;
\item
  absence of distress;
\item
  rapid return to productivity;
\item
  restoration of every pre-event belief;
\item
  forgiveness on demand;
\item
  or proof that harm was beneficial.
\end{itemize}

Repair can include changed boundaries, disability accommodation, legal
action, medication, social recognition, grief, new identity, and refusal
to return to the prior environment.

\subsection{7.9 Post-traumatic growth}\label{post-traumatic-growth}

Reports of growth after adversity may reflect real change, coping
narratives, social expectations, or mixtures. Growth in one dimension
can coexist with symptoms in another. The framework therefore treats
post-traumatic growth as a possible profile change, not a required stage
or retrospective justification of suffering.

\subsection{7.10 Ongoing threat and structural
conditions}\label{ongoing-threat-and-structural-conditions}

A person cannot fully ``regulate'' their way out of continuing violence,
poverty, racism, war, unsafe housing, coercive employment, or abusive
dependence. Clinical methods can support survival and agency, but they
do not remove the environment.

Trauma-informed practice becomes incoherent when it translates political
or material danger into an individual's nervous-system deficit.
Environmental repair and institutional accountability may be primary
psychological interventions.

\subsection{7.11 Falsifiable
implications}\label{falsifiable-implications}

The profile architecture predicts that two people with similar symptom
totals may require different interventions because their binding
constraints differ. It predicts that improvement in distress can occur
without improvement in agency or relationship, and vice versa. It
predicts that the same avoidance behavior can be protective under
ongoing threat and maintaining under current safety.

These are not unique to CU. The framework earns value only if its
declared profile and horizon improve measurement, case formulation, or
prediction over existing approaches.

\begin{center}\rule{0.5\linewidth}{0.5pt}\end{center}

\section{8. Neurodiversity, psychopathology, and
diagnosis}\label{neurodiversity-psychopathology-and-diagnosis}

\subsection{8.1 Difference, disability, distress, and
disorder}\label{difference-disability-distress-and-disorder}

Neurodiversity names the fact and value of variation in cognitive,
sensory, attentional, communicative, and regulatory organization. It
also names a movement challenging research and services organized
exclusively around normalization. It does not imply that every
difficulty is socially produced or that support is unnecessary.

This paper distinguishes:

\begin{itemize}
\tightlist
\item
  \textbf{difference:} variation from a reference distribution;
\item
  \textbf{disability:} restriction in activity or participation arising
  through person--environment interaction;
\item
  \textbf{distress:} suffering reported or behaviorally evident;
\item
  \textbf{impairment:} loss or limitation of a function relevant to the
  person's life;
\item
  \textbf{danger:} risk of serious harm;
\item
  \textbf{disorder:} a clinical classification under a diagnostic
  system; and
\item
  \textbf{injustice:} avoidable harm caused by exclusion, stigma,
  coercion, or unequal resources.
\end{itemize}

These can overlap without being identical.

\subsection{8.2 Beyond ``alternative
coherence''}\label{beyond-alternative-coherence}

Neurodiversity should not be reduced to ``alternative coherence.'' That
phrase risks romanticizing disability and suggesting that every clinical
difficulty is merely social mismatch.

The revised claim is:

\begin{quote}
\textbf{{[}Bridge principle{]} A neurocognitive profile should be
evaluated across person-specific capacities, distress, support needs,
environment fit, strengths, safety, and goals rather than against one
norm of psychological organization.}
\end{quote}

\subsection{8.3 Autism and reciprocal
mismatch}\label{autism-and-reciprocal-mismatch}

Neurodiversity-informed autism research criticizes an exclusive deficit
focus, excessive individualization, and lack of autistic participation
(Pellicano \& den Houting, 2022). The double-empathy hypothesis proposes
that communication difficulty can arise between differently organized
people rather than residing solely in the autistic person.

Evidence is informative and mixed. A 2020 diffusion-chain study found
effective information transfer within autistic groups and reduced
transfer in mixed-neurotype groups (Crompton et al., 2020). A larger
preregistered 2025 Registered Report did not replicate the mixed-chain
transfer deficit, though it found differences in rapport and effects of
diagnostic disclosure (Crompton et al., 2025). The correct conclusion is
not that one result settles the theory. It is that reciprocal fit, task,
heterogeneity, knowledge, and method require explicit treatment.

\subsection{8.4 Camouflaging and
masking}\label{camouflaging-and-masking}

Camouflaging can include suppression, imitation, compensation,
rehearsal, environmental control, and selective disclosure. It may
increase access, safety, employment, or communication while carrying
costs in effort, authenticity, stress, or delayed recognition.

Recent reviews emphasize definitional heterogeneity and sample
limitations (Nel, Spedding, \& Malcolm-Smith, 2025; Greig, Coundouris,
\& Henry, 2026). Cross-sectional associations between masking, distress,
self-esteem, and trauma do not alone establish simple causal direction.

Masking is therefore a clear local--global and person--environment case.
A strategy can be chosen, coerced, protective, skilled, costly, or all
of these.

\subsection{8.5 Clinical heterogeneity}\label{clinical-heterogeneity}

Traditional diagnoses support communication, service access,
epidemiology, and treatment guidelines. They also combine heterogeneous
presentations and produce high comorbidity.

Several complementary approaches address this:

\begin{itemize}
\tightlist
\item
  \textbf{dimensional taxonomies}, including HiTOP;
\item
  \textbf{research-domain approaches}, including RDoC;
\item
  \textbf{general liability models}, including the p factor;
\item
  \textbf{symptom-network models};
\item
  \textbf{process-based formulations}; and
\item
  \textbf{idiographic longitudinal models}.
\end{itemize}

None is complete. HiTOP has substantial support and continues to undergo
extension and validation. Network and process-based methods offer causal
hypotheses but are often applied cross-sectionally or at group level
even when individual dynamics are the target (Borsboom, 2017; Kotov et
al., 2017; Fisher, Medaglia, \& Jeronimus, 2018; Hofmann \& Hayes, 2019;
Schumacher et al., 2024; Westhoff et al., 2024).

CU should not claim that ``coherence'' has already unified these
programs. It can provide a translation architecture for comparing what
each preserves and omits.

\subsection{8.6 Psychopathology without one
essence}\label{psychopathology-without-one-essence}

\textbf{{[}Definition{]} Psychopathology} in this paper refers to
patterns of experience, behavior, or functioning associated with
clinically significant distress, disability, danger, loss of agency, or
impairment, assessed under an explicit diagnostic and cultural context.

A condition can involve:

\begin{itemize}
\tightlist
\item
  distorted prediction;
\item
  maladaptive learning;
\item
  dysregulated affect;
\item
  rigid behavior;
\item
  altered salience;
\item
  cognitive impairment;
\item
  biological disease;
\item
  social exclusion;
\item
  traumatic adaptation;
\item
  disrupted sleep;
\item
  substance effects;
\item
  meaning loss;
\item
  or combinations.
\end{itemize}

The framework rejects both ``all disorder is brain disease'' and ``all
disorder is social construction.'' The causal share differs by condition
and person.

\subsection{8.7 Symptom networks and
feedback}\label{symptom-networks-and-feedback}

Symptoms can causally reinforce one another: insomnia increases fatigue;
fatigue reduces activity; inactivity increases isolation; isolation
increases rumination; rumination disrupts sleep. Network models make
this structure explicit.

But a network estimated from between-person correlations is not
automatically the causal network of one individual. Measurement timing,
omitted variables, common causes, and estimation choices matter.
Intensive longitudinal data and intervention are needed for stronger
claims.

\subsection{8.8 The p factor and scalar
temptation}\label{the-p-factor-and-scalar-temptation}

General psychopathology factors summarize covariance across symptoms.
They can be useful for prediction and classification. They do not prove
that all suffering is one underlying entity, and they should not be
renamed ``global incoherence.'' A statistical factor depends on
measures, sample, model, and level of analysis.

\subsection{8.9 Diagnosis and lived
experience}\label{diagnosis-and-lived-experience}

A diagnosis can provide recognition, language, treatment access,
community, legal protection, or relief from blame. It can also produce
stigma, reification, diagnostic overshadowing, or a narrowed future.

Assessment should include the person's own account, while not making any
one account immune to evidence. Lived-experience participation improves
research and service design when it carries real authority rather than
symbolic consultation.

\subsection{8.10 Culture}\label{culture}

Cultures shape emotion concepts, selfhood, family obligation, spiritual
experience, gender roles, help-seeking, and the threshold at which
distress is named as illness. Cultural variation does not eliminate
reality or suffering. It changes presentation, interpretation,
resources, and acceptable response.

A psychological profile must therefore ask whether an apparent
inconsistency reflects pathology, bilingual or bicultural adaptation,
multiple roles, migration, minority stress, or the reasonable management
of unsafe contexts.

\subsection{8.11 Gender and identity}\label{gender-and-identity}

Masculinity and femininity are not treated as universal psychological
coherence strategies. Gendered behavior is shaped by biology,
development, identity, culture, power, role, and individual variation.
Archetypal yin--yang language may remain meaningful in poetry or
spiritual interpretation; it is not an empirical taxonomy of men and
women.

Distress related to gender can involve bodily incongruence, dysphoria,
stigma, family conflict, discrimination, uncertainty, or combinations.
Technical claims require dedicated evidence and should not be inferred
from framework symbolism.

\subsection{8.12 Disability and
accommodation}\label{disability-and-accommodation}

A profile model must represent the environment. A capacity limitation
may become disabling only under a particular sensory, communicative,
economic, or institutional design. Accommodation can improve functioning
without changing the person.

This is not merely kindness. It is a causal intervention on the system
being modeled.

\begin{center}\rule{0.5\linewidth}{0.5pt}\end{center}

\section{9. Therapy, medication, and other
interventions}\label{therapy-medication-and-other-interventions}

\subsection{9.1 Therapy is not one
mechanism}\label{therapy-is-not-one-mechanism}

Psychotherapy contains many practices with different theories, targets,
and evidence. Some treatments aim at fear learning, behavior, beliefs,
relationships, attention, acceptance, mentalization, emotion regulation,
trauma processing, skills, or meaning. Similar outcomes can arise
through different routes, and similar techniques can work differently
across people.

``Coherence modulation'' may serve as a broad interpretive description
of therapy, but it is not a mechanism.

A scientifically useful formulation must specify:

\begin{itemize}
\tightlist
\item
  the target problem and person-defined goals;
\item
  the proposed mechanism;
\item
  the technique;
\item
  the relevant relationship and context;
\item
  the expected time course;
\item
  adverse effects and contraindications;
\item
  alternative explanations;
\item
  and the outcome profile.
\end{itemize}

\subsection{9.2 Common factors and specific
processes}\label{common-factors-and-specific-processes}

The therapeutic alliance is robustly associated with outcome across
modalities, populations, and delivery formats, though association does
not settle causal direction (Flückiger et al., 2018; Wampold \& Imel,
2015). Hope, expectancy, collaboration, credible rationale, empathy, and
a structured opportunity for new experience may also contribute.

These are enabling conditions, not proof that techniques do not matter.
Trauma-focused exposure, behavioral activation, response prevention,
skills training, interpersonal work, cognitive change, and other
procedures can have disorder- and process-specific effects.

The common-factors debate is therefore misframed when posed as
relationship \emph{versus} technique. A technique is delivered through a
relationship and a relationship becomes therapeutic partly through what
it enables.

\subsection{9.3 Outcomes are multidimensional and often
modest}\label{outcomes-are-multidimensional-and-often-modest}

A 2024 meta-analysis across eight mental disorders found psychotherapies
generally superior to control conditions while absolute response rates
remained modest and varied by condition (Cuijpers et al., 2024). This
matters for claim discipline.

Treatment success should not be reported only as a mean symptom change.
Relevant outcomes may include:

\begin{itemize}
\tightlist
\item
  response and remission;
\item
  functioning;
\item
  quality of life;
\item
  agency;
\item
  relationship;
\item
  sleep;
\item
  safety;
\item
  meaning;
\item
  relapse;
\item
  adverse effects;
\item
  burden and accessibility;
\item
  and the person's own priorities.
\end{itemize}

Nonresponse is not proof that a person resisted coherence. It may
reflect a mismatched formulation, inadequate dose, inaccessible
conditions, ongoing danger, comorbidity, biological factors, clinician
limitations, or an ineffective intervention.

\subsection{9.4 Case formulation}\label{case-formulation}

A CU-informed case formulation is a hypothesis about how processes
interact, not a hidden diagnosis.

It may ask:

\begin{enumerate}
\def\labelenumi{\arabic{enumi}.}
\tightlist
\item
  What does the person want to change or preserve?
\item
  Which symptoms, relationships, environments, or meanings are most
  disabling?
\item
  Which short-term strategies maintain longer-term costs?
\item
  Which strategies remain necessary because threat is ongoing?
\item
  What capacities and supports are available?
\item
  Which protected floors cannot be traded away?
\item
  Which observations would disconfirm the formulation?
\item
  How will the person participate in revising it?
\end{enumerate}

The formulation should become simpler or be abandoned when it adds no
useful prediction.

\subsection{9.5 Cognitive and behavioral
interventions}\label{cognitive-and-behavioral-interventions}

Cognitive and behavioral therapies include diverse methods: exposure,
behavioral activation, cognitive restructuring, response prevention,
skills, experiments, problem solving, and relapse planning.

Their common strength is not that thought alone causes suffering. It is
the use of explicit, testable hypotheses about maintaining processes. A
behavioral experiment can compare predictions. Exposure can create new
learning under tolerable conditions. Behavioral activation can increase
access to reward and structure. Cognitive work can examine
interpretation and standards.

Mechanisms remain debated. Salkovskis, Sighvatsson, and Sigurdsson
(2024) argue that effective therapies help people identify where
maintaining processes have left them stuck and test more flexible
responses. This paper treats this as a useful synthesis rather than a
settled universal mechanism.

\subsection{9.6 Acceptance, mindfulness, and psychological
flexibility}\label{acceptance-mindfulness-and-psychological-flexibility}

Acceptance-based methods alter the relation to thoughts, feelings, and
sensations rather than requiring their elimination. Meta-analytic
evidence supports changes in several flexibility and inflexibility
processes during acceptance and commitment therapy, with links to
distress outcomes and continuing measurement limitations (Macri \&
Rogge, 2024).

Acceptance must not become passive accommodation to abuse, poverty,
discrimination, or preventable illness. It concerns contact with present
experience and limits; action can include protest, boundary, departure,
treatment, or refusal.

Mindfulness practices can support attention, decentering, compassion,
and emotion regulation. They can also produce distressing or
destabilizing experiences in some people. ``More mindfulness'' is not a
universal dose-response law. Practice, teacher competence, trauma
history, context, and expectations matter (Lindahl et al., 2017).

\subsection{9.7 Psychodynamic, narrative, and relational
work}\label{psychodynamic-narrative-and-relational-work}

Psychodynamic and relational therapies examine recurring patterns,
conflict, defense, attachment, transference, and the ways a person's
history enters present relationships. Narrative therapies examine how
problem-saturated stories organize possibility and how alternative
descriptions can support agency.

This paper translates these approaches into meaning-map and relational
terms without claiming that one vocabulary captures their full practice.
The therapeutic relationship can make previously intolerable affect or
conflict thinkable, but ``borrowed coherence'' is metaphor unless the
regulatory process is operationalized.

\subsection{9.8 Family and systems
interventions}\label{family-and-systems-interventions}

A person's difficulties may be maintained by interaction patterns, role
expectations, caregiving burden, economic stress, or institutional
response. Family and systems interventions can change contingencies and
communication that individual therapy cannot.

Systems language must not erase responsibility. Abuse is not merely a
circular communication pattern. Power, coercion, and safety require
explicit representation.

\subsection{9.9 Trauma-focused
treatment}\label{trauma-focused-treatment}

Clinical guidelines continue to support several individual
trauma-focused psychotherapies for PTSD. Evidence for complex
presentations is expanding while certainty varies and personalization
remains important (VA/DoD, 2023; Billings \& Nicholls, 2025; Schaug et
al., 2025).

Two universal claims should be rejected:

\begin{itemize}
\tightlist
\item
  everyone must complete a fixed stabilization phase before trauma
  processing; and
\item
  immediate confrontation is always superior to staged work.
\end{itemize}

The relevant decision concerns safety, dissociation, skills, preference,
comorbidity, ongoing threat, treatment access, and the evidence for the
specific method.

\subsection{9.10 Somatic and embodied
approaches}\label{somatic-and-embodied-approaches}

Movement, breathing, posture, interoceptive training, body-oriented
psychotherapy, yoga, and other embodied practices may affect arousal,
attention, pain, sleep, agency, and bodily safety. Evidence quality and
mechanisms vary. A 2024 meta-analysis of yoga for PTSD reported
improvements in self-reported symptoms and depression but not
clinician-rated PTSD, with risk-of-bias limitations in the underlying
trials (Nejadghaderi et al., 2024).

The publication-safe position is:

\begin{quote}
\textbf{{[}Conditional result{]} Embodied practices can be useful
components of care for some people and outcomes. They should not be
presented as universally superior, as proof that trauma is literally
stored in tissue, or as substitutes for indicated treatment.}
\end{quote}

\subsection{9.11 Medication}\label{medication}

Psychiatric medications can alter mood, arousal, attention, sleep,
psychosis, compulsivity, pain, and other processes. Their effects are
heterogeneous, probabilistic, and person-specific. They can support
agency and therapy, produce adverse effects, lose effectiveness, or be
difficult to discontinue.

Simple labels such as the following should be rejected:

\begin{itemize}
\tightlist
\item
  dopamine equals motivation;
\item
  serotonin equals baseline coherence;
\item
  norepinephrine equals threat;
\item
  excitation--inhibition balance equals mental health.
\end{itemize}

Neurotransmitters act through multiple receptors, circuits, timescales,
and developmental contexts. Psychological meaning and social conditions
affect how changes are experienced.

A profile-based medication assessment records target symptoms,
functioning, side effects, sleep, agency, relationships, adherence
burden, withdrawal or discontinuation effects, and person preference. It
does not ask whether the drug ``raises coherence.''

\subsection{9.12 Psychedelic-assisted
therapy}\label{psychedelic-assisted-therapy}

Psychedelic compounds can alter perception, self-processing, affect,
suggestibility, and meaning. Clinical trials report promising effects
for selected conditions and carefully screened participants, while
blinding, expectancy, therapist effects, durability, adverse events, and
generalizability remain important limitations.

REBUS, entropic-brain, predictive-processing, and criticality accounts
are model families, not settled mechanisms (Carhart-Harris \& Friston,
2019). The phrase ``constraint relaxation'' can be useful if the
constraint, measurement, and alternative explanation are declared.

Psychedelic experience does not automatically produce truth, healing, or
moral insight. It may intensify commitment, confusion, grandiosity,
fear, or dependence on an interpretive authority. Preparation, consent,
setting, integration, medical screening, and power are therefore part of
the causal system.

\subsection{9.13 Exercise, sleep, and ordinary
practices}\label{exercise-sleep-and-ordinary-practices}

Exercise, sleep regularity, social connection, time outdoors, nutrition,
creative activity, and structured routine can support mental health.
Their average effects do not make them adequate treatment for every
condition or morally obligatory proof of effort.

Ordinary practices matter partly because they alter biological capacity,
environmental predictability, reward, social contact, and agency. Their
limits matter because a person may lack safety, time, money, mobility,
access, or energy.

\subsection{9.14 Spiritual and contemplative
practices}\label{spiritual-and-contemplative-practices}

Prayer, ritual, chanting, service, confession, pilgrimage, study,
community, and meditation can regulate attention, affect, identity,
belonging, grief, and commitment. They can also intensify shame,
obedience, exclusion, bypass, or spiritual struggle.

The descriptive question asks what a practice does. The commitment
question asks what form of life it enacts. The metaphysical question
asks what, if anything, it discloses about reality. This paper keeps all
three visible.

\subsection{9.15 Measurement-based and feedback-informed
care}\label{measurement-based-and-feedback-informed-care}

Repeated measurement can detect deterioration, nonresponse, and
mismatch. Yet a questionnaire can narrow treatment to what it measures.
Feedback systems should therefore combine symptom and functioning
measures with adverse effects, goals, alliance, context, and qualitative
report.

A profile is useful only when it improves shared decision-making rather
than creating another authority over the person.

\subsection{9.16 The therapeutic stance}\label{the-therapeutic-stance}

The therapeutic stance appropriate to this framework is neither detached
expertise nor unconditional affirmation. It combines:

\begin{itemize}
\tightlist
\item
  respect for the person's experience;
\item
  openness to correction;
\item
  attention to safety and power;
\item
  curiosity about function;
\item
  willingness to test hypotheses;
\item
  protection of agency;
\item
  and honesty about uncertainty and limits.
\end{itemize}

Compassion is epistemic because it helps keep another person's reality
in view. It is not a substitute for evidence or boundary.

\begin{center}\rule{0.5\linewidth}{0.5pt}\end{center}

\section{10. Development, plasticity, and psychological
trajectories}\label{development-plasticity-and-psychological-trajectories}

\subsection{10.1 Development is not one
ladder}\label{development-is-not-one-ladder}

Psychological development involves changing relations among embodiment,
attachment, language, attention, emotion regulation, identity, social
role, moral reasoning, and agency. These capacities do not mature in one
order or at one rate.

Development should not be represented as one universal sequence
culminating in a single spiritual-ethical profile. Culturally variable
forms of life, asynchronous capacities, and different biographies
require a \textbf{developmental profile} and a model family of path
dependence.

\subsection{10.2 Developmental profile}\label{developmental-profile}

A developmental description may track:

\begin{itemize}
\tightlist
\item
  physiological regulation;
\item
  attachment and trust;
\item
  language and symbolic capacity;
\item
  executive function;
\item
  emotion understanding and regulation;
\item
  play and exploration;
\item
  self-recognition;
\item
  autonomy;
\item
  peer belonging;
\item
  identity;
\item
  moral and epistemic agency;
\item
  intimacy;
\item
  caregiving;
\item
  generativity;
\item
  and adaptation to aging, illness, and loss.
\end{itemize}

These dimensions are shaped by culture and opportunity. A person may
show advanced reasoning with limited emotional regulation, strong care
with limited autonomy, or deep spiritual commitment with unresolved
trauma. No single total rank follows.

\subsection{10.3 External regulation and internal
capacity}\label{external-regulation-and-internal-capacity}

Early development depends heavily on caregivers and environments.
Co-regulation helps organize arousal, sleep, attention, exploration, and
expectations of response. Over time, people may internalize skills,
routines, language, and self-soothing while continuing to depend on
others.

Development does not move from dependence to independence in a simple
line. Mature functioning is often interdependent: people know when to
act alone, when to seek help, and how to contribute to shared
regulation.

\subsection{10.4 Path dependence}\label{path-dependence}

Repeated interactions can stabilize patterns that resist small
perturbations. Attractor language can describe this qualitatively, but a
mathematical attractor requires declared state variables and dynamics.

\textbf{{[}Bridge principle{]} Development is path-dependent when
earlier states and environments alter the options, costs, or transitions
available later.}

Path dependence does not imply irreversibility. Plasticity can persist,
and new relationships, environments, skills, medication, injury,
practice, or social recognition can reorganize trajectories.

\subsection{10.5 Sensitive periods and ongoing
plasticity}\label{sensitive-periods-and-ongoing-plasticity}

Some capacities are easier to acquire during particular developmental
windows. This does not mean that later change is impossible or that one
missed experience determines destiny. Claims about sensitive periods
must be domain-specific.

The language of ``rewiring'' is often too broad. Change can involve new
learning, inhibition of old responses, environmental scaffolding,
narrative revision, bodily regulation, altered incentives, or structural
support without erasing prior patterns.

\subsection{10.6 Adolescence and
identity}\label{adolescence-and-identity}

Adolescence often intensifies peer evaluation, autonomy, future
planning, sexuality, and identity exploration. It can also bring
increased vulnerability to mood, anxiety, substance use, and social
exclusion.

A profile approach avoids treating disagreement with adults as pathology
or every risk as healthy exploration. It asks which capacities,
constraints, and supports are changing and whose goals define success.

\subsection{10.7 Adulthood}\label{adulthood}

Adult development continues through work, relationship, caregiving,
migration, illness, parenthood, loss, moral conflict, and changing
embodiment. Trait stability coexists with change. Narrative identity can
reorganize in response to major events.

Competence in one role can conceal fragility in another. Productivity is
not a complete measure of psychological health. Nor is subjective
comfort when it depends on exploitation, intoxication, or systematic
self-deception.

\subsection{10.8 Aging}\label{aging}

Aging changes sensory function, energy, memory, social role, bodily
capacity, and exposure to loss. It can also support emotional
regulation, perspective, selective investment, and meaning.

Aging is not a universal ``coherence decline.'' Different functions
change differently. Disability and dependence do not erase personhood,
value, agency, or meaning.

\subsection{10.9 Maturity}\label{maturity}

Psychological maturity is a normative and empirical hybrid. Candidate
capacities include:

\begin{itemize}
\tightlist
\item
  reality contact;
\item
  affect tolerance;
\item
  flexible strategy use;
\item
  responsibility;
\item
  reciprocal relationship;
\item
  capacity for repair;
\item
  stable but revisable identity;
\item
  and action across longer horizons.
\end{itemize}

No evidence establishes one final profile for all cultures and lives.
Truth joined with compassion remains an important CU commitment, but it
is not a developmental theorem.

\subsection{10.10 Regression, recurrence, and
context}\label{regression-recurrence-and-context}

A capacity available in one context may disappear under threat, sleep
loss, illness, intoxication, grief, coercion, or overwhelming demand.
This is not necessarily a return to an earlier developmental stage. It
may be a state-dependent reduction in access.

Assessment should distinguish loss of capacity from loss of performance
under current conditions.

\begin{center}\rule{0.5\linewidth}{0.5pt}\end{center}

\section{11. Psychological infrastructure and the meaning-crisis
hypothesis}\label{psychological-infrastructure-and-the-meaning-crisis-hypothesis}

\subsection{11.1 Minds are not
self-contained}\label{minds-are-not-self-contained}

Individual regulation depends on housing, income, food, healthcare,
education, transport, community, discrimination, work, family,
technology, ecology, and political stability. The WHO's 2022 report and
systematic reviews of social determinants emphasize that mental health
is shaped by environments and systems, not only by individual traits
(Alegría et al., 2018; Lund et al., 2018; World Health Organization,
2022).

This paper calls these conditions \textbf{psychological coherence
infrastructure} when they support reality orientation, agency,
belonging, regulation, development, and repair. The term is a bridge,
not a replacement for social epidemiology.

\subsection{11.2 Infrastructure profile}\label{infrastructure-profile}

Relevant supports may include:

\begin{itemize}
\tightlist
\item
  material security;
\item
  bodily safety;
\item
  time and rest;
\item
  access to care;
\item
  reliable information;
\item
  relational belonging;
\item
  recognition and rights;
\item
  meaningful participation;
\item
  cultural continuity;
\item
  green and public space;
\item
  and institutions capable of correction.
\end{itemize}

A person can possess strong internal capacities and still be overwhelmed
by sustained external demand. A system can also provide material support
while undermining dignity, truth, or agency.

\subsection{11.3 The meaning crisis as a hypothesis
family}\label{the-meaning-crisis-as-a-hypothesis-family}

``Meaning crisis'' can refer to several distinct claims:

\begin{itemize}
\tightlist
\item
  reduced experienced purpose;
\item
  reduced comprehensibility;
\item
  reduced mattering;
\item
  loneliness or weak belonging;
\item
  declining institutional trust;
\item
  loss of shared narrative;
\item
  economic precarity;
\item
  ecological threat;
\item
  religious change;
\item
  attention capture;
\item
  or conflict among inherited values and present conditions.
\end{itemize}

These trends do not necessarily covary. A society may lose one shared
narrative while expanding plural forms of meaning. A person may reject
inherited religion and gain purpose. Another may gain information access
and lose trust.

The framework should therefore speak of \textbf{meaning-crisis
hypotheses}, each with measures, population, period, comparison, and
mechanism.

\subsection{11.4 Digital environments}\label{digital-environments}

Digital systems can provide knowledge, relationship, creativity,
accessibility, and community. They can also amplify comparison, outrage,
misinformation, surveillance, compulsive use, and attentional
fragmentation.

``Technology reduces coherence'' is too broad. The relevant variables
include platform design, incentives, content, user goals, social
context, age, vulnerability, and available alternatives.

\subsection{11.5 Work and economic life}\label{work-and-economic-life}

Employment can provide income, routine, identity, competence, social
connection, and contribution. It can also produce precarity, injury,
alienation, surveillance, moral conflict, and exhaustion.

Psychological language becomes ideological when it teaches people to
tolerate conditions that should be changed. Resilience programs cannot
substitute for fair workload, safety, autonomy, and material security.

\subsection{11.6 Institutions and epistemic
support}\label{institutions-and-epistemic-support}

People depend on institutions for records, expertise, adjudication,
education, and correction. When institutions become unreliable or are
perceived as illegitimate, individuals face greater burdens of
verification and distrust.

Trust is not blind deference. Healthy epistemic infrastructure supports
criticism, transparency, accountability, and revision.

\subsection{11.7 Religion and shared
meaning}\label{religion-and-shared-meaning}

Religious communities can provide belonging, ritual, intergenerational
memory, moral language, service, and ways of holding suffering. They can
also enforce conformity, shame, exclusion, abuse, or factual insulation.

Calling religion a ``coherence technology'' is too instrumental.
Traditions are forms of life whose meanings exceed their psychological
functions. Psychology can study their effects without reducing devotion
to regulation.

\subsection{11.8 Polarization and identity
threat}\label{polarization-and-identity-threat}

Political and cultural conflict can become psychologically
self-reinforcing when identity, belonging, status, fear, and information
sources align. Opposing evidence then threatens more than a belief; it
threatens the relational and moral world in which the person lives.

This helps explain resistance without excusing falsehood or harm. The
intervention problem includes material incentives, group norms, trusted
messengers, information architecture, and opportunities for cross-group
recognition---not merely individual debiasing.

\subsection{11.9 Collective trauma and
memory}\label{collective-trauma-and-memory}

Communities can carry histories of war, enslavement, colonization,
displacement, genocide, repression, and institutional betrayal through
narratives, practices, bodies, laws, and unequal conditions.
``Collective trauma'' should not be treated as one group mind or
inherited memory substance.

The term is useful when mechanisms of transmission are specified: family
interaction, education, silence, discrimination, material deprivation,
commemorative practice, and continuing threat.

\subsection{11.10 Implications for care}\label{implications-for-care}

A clinically responsible framework asks at three levels:

\begin{enumerate}
\def\labelenumi{\arabic{enumi}.}
\tightlist
\item
  What can this person learn or change?
\item
  What relationship or community support is needed?
\item
  What environmental or institutional condition must be altered?
\end{enumerate}

Focusing only on the first can become victim-blaming. Focusing only on
the third can neglect capacities that matter while structural change is
incomplete.

\begin{center}\rule{0.5\linewidth}{0.5pt}\end{center}

\section{12. Empirical program, measurement, and
falsification}\label{empirical-program-measurement-and-falsification}

\subsection{12.1 The research object}\label{the-research-object}

The framework's central empirical claim is not that all psychology
reduces to coherence. It is that declared profiles, horizons, and
relational models can improve explanation where isolated variables or
undifferentiated totals fail.

The research program therefore compares models.

\subsection{12.2 Measurement program}\label{measurement-program}

A psychological coherence profile should be assembled from validated
measures where available rather than invented as one new omnibus
questionnaire. Candidate components can be measured through:

\begin{itemize}
\tightlist
\item
  calibration and reality-testing tasks;
\item
  emotion-regulation measures;
\item
  psychological flexibility measures;
\item
  narrative coding;
\item
  attachment and relational functioning;
\item
  behavioral agency;
\item
  sleep and physiological measures;
\item
  meaning-in-life measures;
\item
  functioning and disability;
\item
  and repair or recovery trajectories.
\end{itemize}

Measurement invariance across culture, language, disability, age, and
context must be tested. A profile built from instruments that encode one
normative population can reproduce that norm under technical language.

\subsection{12.3 Intensive longitudinal and idiographic
models}\label{intensive-longitudinal-and-idiographic-models}

Group averages do not automatically generalize to one person. Intensive
longitudinal assessment can estimate temporal relations among sleep,
affect, conflict, avoidance, pain, activity, and support. Idiographic
models are especially relevant when the framework claims that binding
constraints differ across people.

But dense data do not guarantee causality. Reactivity, missingness,
timing, confounding, and model instability remain.

\subsection{12.4 Intervention and
ablation}\label{intervention-and-ablation}

A proposed maintaining process gains credibility when targeted change
produces predicted effects while plausible alternatives are measured.
Examples include:

\begin{itemize}
\tightlist
\item
  reducing avoidance under safe conditions;
\item
  increasing sleep opportunity;
\item
  altering social contingencies;
\item
  changing a belief through experiment;
\item
  adding accommodation;
\item
  or changing medication.
\end{itemize}

Intervention must be ethical and person-centered. The framework does not
license experimental destabilization merely because a pattern appears
rigid.

\subsection{12.5 Profile predictions}\label{profile-predictions}

The framework generates several testable predictions:

\begin{enumerate}
\def\labelenumi{\arabic{enumi}.}
\tightlist
\item
  Profile models will distinguish people with similar symptom totals but
  different treatment needs.
\item
  Immediate relief and long-horizon agency will sometimes dissociate.
\item
  A strategy's effect will change with environmental threat and support.
\item
  Narrative clarity will not always predict descriptive adequacy or
  flexibility.
\item
  Changes in alliance, skill, behavior, meaning, and symptoms will show
  different temporal sequences across persons.
\item
  Accommodation can improve functioning without normalizing the
  underlying cognitive profile.
\item
  Self-deception judgments based on convergent asymmetry evidence will
  be more reliable than judgments based on false belief or disagreement
  alone.
\item
  Social and material interventions will explain variance that
  individual-only models classify as residual vulnerability.
\end{enumerate}

\subsection{12.6 Competing explanations}\label{competing-explanations}

For every result, the framework should compare at least one non-CU
model. Possible competitors include:

\begin{itemize}
\tightlist
\item
  symptom severity alone;
\item
  diagnostic category;
\item
  trait models;
\item
  reinforcement learning;
\item
  predictive processing;
\item
  attachment;
\item
  network models;
\item
  social determinants;
\item
  and simple resource limitation.
\end{itemize}

A CU translation is useful only if it clarifies relations among these
models or generates discriminating predictions.

\subsection{12.7 Falsification and failure
conditions}\label{falsification-and-failure-conditions}

The framework should be revised or abandoned in a domain if:

\begin{itemize}
\tightlist
\item
  profile components cannot be measured with acceptable reliability;
\item
  profiles add no predictive or clinical value over established
  constructs;
\item
  local--global or horizon analyses do not improve decision quality;
\item
  meaning maps cannot be operationalized beyond metaphor;
\item
  self-deception profiles show poor inter-rater reliability or
  systematic ideological bias;
\item
  infrastructure variables do not contribute after appropriate controls;
\item
  claimed cross-domain translations collapse into equivocation;
\item
  or the framework repeatedly normalizes culturally preferred lives
  while calling the result coherence.
\end{itemize}

\subsection{12.8 Ethical constraints on the research
program}\label{ethical-constraints-on-the-research-program}

Psychological research involves persons whose reports, identities, and
vulnerabilities can be affected by measurement. Governance should
include:

\begin{itemize}
\tightlist
\item
  informed consent;
\item
  privacy;
\item
  data minimization;
\item
  lived-experience participation;
\item
  cultural competence;
\item
  adverse-event monitoring;
\item
  protection against coercive use;
\item
  and limits on automated diagnosis.
\end{itemize}

A profile system can be used to support care or to rank, surveil, and
exclude. The architecture does not determine the ethics of its
deployment.

\subsection{12.9 Relationship to clinical
practice}\label{relationship-to-clinical-practice}

This paper presents a research architecture. It does not create a CU
diagnosis, certify therapists, or authorize treatment claims. Clinical
adoption would require validated instruments, trials, training,
regulation, and evidence of benefit beyond existing approaches.

\subsection{12.10 What would count as
success}\label{what-would-count-as-success}

Success would not be universal adoption of the vocabulary. It would be a
smaller achievement:

\begin{itemize}
\tightlist
\item
  better distinctions;
\item
  fewer category errors;
\item
  more transparent assumptions;
\item
  improved person-specific prediction;
\item
  more honest representation of tradeoffs;
\item
  and interventions that preserve agency while reducing suffering.
\end{itemize}

\begin{center}\rule{0.5\linewidth}{0.5pt}\end{center}

\begin{center}\rule{0.5\linewidth}{0.5pt}\end{center}

\section{13. Psychology, Healing, and Sacred
Repair}\label{psychology-healing-and-sacred-repair}

\subsection{13.1 Register declaration}\label{register-declaration}

This section is interpretive. It is constrained by the psychological
account but not entailed by it. It does not make a clinical claim,
recommend treatment for a particular person, or assign spiritual meaning
to another person's suffering without their consent.

Psychological science can investigate memory, threat, attachment,
affect, behavior, embodiment, and change. The systems-metaphysical
interpretation places a human life within the wider history of living
and conscious organization. The commitment register asks how that life
should be met.

\subsection{13.2 Healing as expanded truthful
capacity}\label{healing-as-expanded-truthful-capacity}

Healing is not seamless consistency or the elimination of pain.

\textbf{{[}Interpretation{]}} Healing is the widening of what a life can
truthfully feel, know, remember, relate to, and act from without
abandoning itself or erasing others.

This widening may involve greater access to previously excluded
experience, but access is not always the first need. Safety, distance,
medication, rest, anger, numbness, boundary, material support, or
temporary compartmentalization may be necessary conditions of later
integration.

Repair is therefore timed. What would be avoidance in one context may be
protection in another. What once protected life may later restrict it.
Healing requires enough truth to see both.

\subsection{13.3 Trauma as contraction of
possibility}\label{trauma-as-contraction-of-possibility}

Trauma can narrow the range of states, memories, relationships, bodily
sensations, meanings, and futures that can be approached without
overwhelming threat.

The contraction is not a moral failure. It can be the cost of survival
under conditions in which openness was unsafe. Hypervigilance,
avoidance, constriction, dissociation, appeasement, aggression, and
compartmentalization may all carry protective logic.

Sacred repair begins by honoring that logic without making it permanent.
It asks what the response protected, what it now costs, what conditions
have changed, and what support would allow new possibilities to become
tolerable.

\subsection{13.4 The body is not beneath
meaning}\label{the-body-is-not-beneath-meaning}

Psychological life is embodied. A person may understand a truth
conceptually while remaining unable to enact it under threat. Bodily
regulation, sleep, pain, movement, medication, sensory conditions,
nutrition, illness, disability, and safety can determine which
psychological capacities are available.

The interpretive commitment of Heaven≡Earth refuses to treat insight as
superior to embodiment. A spiritual or therapeutic claim that cannot
land in the body's actual conditions remains incomplete.

\subsection{13.5 Relationship and borrowed
possibility}\label{relationship-and-borrowed-possibility}

Human regulation is partly distributed. A trusted relationship can
change perceived effort, threat, attention, memory, and access to
action. Care can lend a person capacities they cannot yet carry alone.

This does not make dependency a defect. Nor does it make every
relationship safe. Sacred relation requires consent, boundary,
reciprocity appropriate to role, and the ability to refuse or leave
where possible.

Healing sometimes occurs because another person can remain present
without demanding that pain become intelligible too quickly.

\subsection{13.6 Ego, identity, and non-egoic
life}\label{ego-identity-and-non-egoic-life}

The ego is not an enemy substance. It is a family of self-modeling and
self-protective processes that can preserve continuity, boundary, and
action. It becomes imprisoning when the image of self must be protected
from every loss, correction, dependency, or change.

Non-egoic development is not self-annihilation. It is the loosening of
compulsory defense around identity, allowing responsibility without
total shame, conviction without certainty, and care without possession.

\subsection{13.7 Grief, love, and continuing
bonds}\label{grief-love-and-continuing-bonds}

Grief reveals that identity is relational. The dead, the lost, and the
changed continue through memory, practice, bodily expectation, altered
possibility, and love.

Healing need not mean severing the bond. It may mean allowing the bond
to change form while refusing denial of the loss. A worldview that
forbids grief in order to protect certainty has confused spiritual
commitment with emotional insulation.

\subsection{13.8 Spiritual practice and psychological
change}\label{spiritual-practice-and-psychological-change}

Prayer can orient attention and trust. Ritual can hold grief and
continuity. Meditation can loosen identification. Service can transform
meaning into action. Movement can restore embodied possibility.
Community can lend regulation and purpose.

Their spiritual meanings exceed measured psychological effects, but
their uses remain answerable to those effects. A practice that
repeatedly produces dissociation, coercion, dependency, shame, or
avoidance of necessary care cannot claim immunity because its language
is sacred.

\subsection{13.9 Heaven and Earth}\label{heaven-and-earth}

Within the project's symbolic vocabulary:

\begin{itemize}
\tightlist
\item
  \textbf{Heaven} names possibility, abstraction, interpretation, and
  orienting meaning;
\item
  \textbf{Earth} names body, consequence, limitation, relationship, and
  the conditions under which meaning is lived.
\end{itemize}

Heaven≡Earth is the commitment not to sever them. Meaning must become
embodied without being reduced to mechanism. Embodiment must be honored
without treating measurement as the whole of a life.

\subsection{13.10 Sacred healing never sanctifies
harm}\label{sacred-healing-never-sanctifies-harm}

Suffering is not automatically a lesson, punishment, purification, gift,
or necessary stage in awakening. Some suffering enlarges a life through
work done in response. Some merely injures. Some is imposed by injustice
and should have been prevented.

\textbf{{[}Commitment{]}} The meaning made after harm belongs to the
person and community doing the work. It does not retroactively justify
the violence, betrayal, deprivation, illness, or loss that made healing
necessary.

\begin{center}\rule{0.5\linewidth}{0.5pt}\end{center}

\subsection{13.11 Meaning can move}\label{meaning-can-move}

\textbf{{[}Commitment and psychological interpretation{]}}
\textbf{CU-COM16 --- Freedom of Meaning-Revision} names a constrained
but consequential freedom. The event does not change. The injury is not
redescribed out of existence. But persons and communities can sometimes
revise what an event means, how it is integrated, and which futures it
permits.

This freedom is never merely private will. It depends on bodily safety,
time, material stability, language, relationship, memory, therapeutic or
contemplative support, and access to real alternatives. A demand to
``reframe'' injustice while leaving harmful conditions intact is not
meaning-revision; it is abandonment disguised as optimism.

\subsection{13.12 Redemptive possibility without
justification}\label{redemptive-possibility-without-justification}

\textbf{{[}Commitment{]}} \textbf{CU-COM17 --- Redemptive Possibility
Without Justification of Evil} preserves the possibility that suffering
can become solidarity, wisdom, creation, care, justice, or repair.
Redemption changes what suffering becomes in a life or history. It never
makes the original harm good, necessary, deserved, or justified.

No person owes redemption. No victim owes forgiveness, testimony,
growth, reconciliation, or a spiritually satisfying ending. Some losses
remain irreversible and some suffering remains without adequate meaning.
The commitment is to keep open the possibility that the perpetrator,
event, or wound need not possess final authority over the future.

\subsection{13.13 From Omega glimpse to
embodiment}\label{from-omega-glimpse-to-embodiment}

\textbf{{[}Interpretation and research architecture{]}} The Immanent
Omega can be approached psychologically through four distinctions: a
\textbf{glimpse} of widened love or self-transcendence;
\textbf{stabilization} of that orientation across ordinary conditions;
\textbf{embodiment} in truthful, non-dominating conduct; and an
\textbf{ecology} that makes such lives more possible.

Peak intensity is not psychological integration. Romantic merger,
collective effervescence, contemplative absorption, or psychedelic unity
may coexist with grandiosity, avoidance, exploitation, or fragile
functioning. The relevant test is whether the experience expands the
capacity for truth, preserved differentiation, compassion, agency,
accountability, and repair.

\section{14. What Commitments This Account
Licenses}\label{what-commitments-this-account-licenses}

\begin{quote}
\textbf{Register note.} ``Licenses'' does not mean ``clinically proves''
or ``spiritually guarantees.'' The commitments below are orientations
made intelligible by the psychological account. They remain constrained
by evidence, safety, affected persons' agency, material conditions, and
the requirement that no meaning-making practice sanctify harm.
\end{quote}

\begin{quote}
\textbf{Official commitment principle --- CU-COM10: Healing as Sacred
Repair.}

Coherence Universalism commits to treating psychological healing as
sacred work: the restoration or widening of a person's capacity to
encounter reality, inhabit the body, bear feeling, remember, choose,
relate, and imagine a future beyond what once had to be survived.

\textbf{Boundary:} Sacred healing never sanctifies harm, forces
disclosure or integration, requires forgiveness or reconciliation, or
replaces evidence-based medical and clinical care.

\textbf{Revision condition:} The commitment must be revised when its
language pressures survivors, pathologizes protective difference, erases
structural conditions, or privileges spiritual meaning over safety and
autonomy.
\end{quote}

\subsection{14.1 The descriptive ground}\label{the-descriptive-ground}

Psychological lives are history-dependent organizations of body,
perception, memory, affect, meaning, identity, relationship, and action.
Protection and repair are context-sensitive. Immediate relief can
conflict with later agency. Integration can heal and can be forced.
Separation can imprison and can preserve safety. No universal score
resolves these tensions.

\subsection{14.2 A life as reality knowing its own
history}\label{a-life-as-reality-knowing-its-own-history}

\textbf{{[}Interpretation{]}} A human being is one place where reality
becomes capable of remembering, anticipating, grieving, interpreting,
and revising its own local history.

Trauma can contract that capacity. Healing can restore the ability to
perceive, feel, remember, choose, and relate without being completely
organized around what once had to be survived.

\begin{quote}
\textbf{Healing is one way the universe comes to know itself more
clearly.}
\end{quote}

This does not imply one cosmic person or make a wounded life valuable
only because it serves a universal process. The person is not an
instrument of the metaphor. The metaphor is answerable to the person.

\subsection{14.3 Psychological healing is sacred
work}\label{psychological-healing-is-sacred-work}

\textbf{{[}Commitment{]} Psychological healing is sacred work.}

It is sacred because a person possesses significance beyond their
productivity, social convenience, diagnostic category, or usefulness to
others. Repair restores access to truth, body, feeling, agency,
relation, imagination, and future possibility.

Sacred work may be quiet. It may look like sleeping, taking medication,
learning a boundary, telling the truth, leaving an unsafe relationship,
grieving, asking for help, tolerating one sensation, or discovering that
survival no longer requires the same defense.

\subsection{14.4 Truth and compassion must remain
joined}\label{truth-and-compassion-must-remain-joined}

Truth without compassion can become abstraction, exposure, or
punishment. Compassion without truth can become collusion, avoidance, or
sentimental control.

\textbf{{[}Commitment{]}} We seek psychological understanding that lets
reality correct us while refusing to reduce a person to the most
measurable or damaged part of their life.

This includes truth about inner defenses and truth about external
conditions. A person should not be asked to regulate themselves into
tolerating what remains violent, exploitative, or unsafe.

\subsection{14.5 Boundaries can be
sacred}\label{boundaries-can-be-sacred}

Integration is not always the immediate good. A boundary can preserve
the conditions under which later contact becomes possible. Distance can
be truthful. Refusal can be an act of agency. Protective
compartmentalization can be the form coherence takes under danger.

\textbf{{[}Commitment{]}} No spiritual or therapeutic ideal of unity may
override consent, timing, safety, or the right not to reconcile.

\subsection{14.6 Meaning without retrospective
justification}\label{meaning-without-retrospective-justification}

A person may create meaning after suffering. They may become more
compassionate, truthful, courageous, or capable. None of this makes the
suffering deserved or necessary.

\textbf{{[}Commitment{]}} We honor growth without converting it into a
defense of what caused the wound.

\subsection{14.7 Meaning, revision, and
redemption}\label{meaning-revision-and-redemption}

\textbf{CU-COM15 --- Meaning as Discovery and Participation} commits CU
to meaning that is objective enough to constrain, relational enough to
be plural, and historical enough to be revised. Meaning is discovered
through truthful participation in real goods, relationships, histories,
responsibilities, and possibilities; it is actualized through
commitment.

\textbf{CU-COM16} protects the freedom through which meaning can move,
while \textbf{CU-COM17} protects redemptive possibility without
converting suffering into destiny or theodicy. Together they make
healing more than symptom reduction while refusing every demand that a
wounded person produce a morally pleasing story.

\subsection{14.8 Love and Omega as ordinary
practice}\label{love-and-omega-as-ordinary-practice}

\textbf{CU-COM14 --- Love as the Ground and Aim of Conscious Life} and
\textbf{CU-COM19 --- Enact Omega Here} orient psychological growth
toward widened, non-possessive concern rather than toward
invulnerability or spiritual status. The point is not to remain in a
peak state. It is to make love, sympathetic joy, truthfulness, service,
boundary, and repair more available in ordinary life.

Omega embodiment is therefore tested by conduct. A person is not
psychologically or spiritually authorized by a mystical experience, a
therapeutic breakthrough, or an identity as ``healed.'' Greater
integration should increase humility, reality contact, capacity for
relationship, and willingness to repair---not exemption from criticism.

\subsection{14.9 The vow of sacred
repair}\label{the-vow-of-sacred-repair}

The commitments of this paper can be gathered as follows:

\begin{quote}
We commit to meet suffering as real.\\
We commit to protect agency and the right to define one's own pace of
repair.\\
We commit to truth about both inner defenses and outer conditions.\\
We commit not to force disclosure, forgiveness, reconciliation, or
integration.\\
We commit to honor the body, relationship, material support, and
spiritual meaning without letting any one erase the others.\\
We commit to build conditions in which lives can widen beyond what they
once had to survive.
\end{quote}

\subsection{14.10 Boundaries of the
commitment}\label{boundaries-of-the-commitment}

This account does not license:

\begin{itemize}
\tightlist
\item
  diagnosing a particular person from CU concepts;
\item
  treating distress as proof of hidden growth;
\item
  treating symptom relief as failure because it is not transformative;
\item
  replacing medical or clinical care with spiritual interpretation;
\item
  demanding reconciliation with perpetrators or unsafe systems;
\item
  declaring every fragmentation pathological;
\item
  or using ``the universe healing itself'' to erase the concrete person
  whose life is at stake.
\end{itemize}

\begin{center}\rule{0.5\linewidth}{0.5pt}\end{center}

\section{15. Conclusion: Truthful Life Across
Change}\label{conclusion-truthful-life-across-change}

Psychological coherence is not seamless consistency. It is the revisable
capacity of a meaning-bearing life to remain in contact with reality,
experience, relationship, body, and action across change---sometimes by
preserving continuity, sometimes by tolerating rupture, and sometimes by
becoming someone who can hold what the former self could not.

The scientific contribution of this paper is architectural. It replaces
one hidden psychological quantity with declared profiles, histories,
constraints, protected floors, observation channels, and competing
explanations. It treats meaning as relational, motivation as plural,
trauma as heterogeneous, protection as context-sensitive, diagnosis as
useful but incomplete, and therapy as a family of evidence-bearing
processes.

The spiritual contribution is not a new clinical mechanism. It is a way
of receiving repair.

A wounded life is not defective matter. It is a living center whose
possibilities have been narrowed by what happened and by what remains.
Healing is sacred when it restores truthful contact and viable freedom
without denying the defenses that made survival possible.

The movement into Ethics now becomes visible. Epistemology asks us not
to lie about reality. Psychology reveals the places where fear, shame,
attachment, and identity can make truth difficult to bear. Ethics asks
what responsibility follows when those same forces shape how we treat
others and justify power.

The task is not to become invulnerable. It is to become more capable of
truth, feeling, relationship, and repair without abandoning the body and
history through which that capacity must live.

\begin{center}\rule{0.5\linewidth}{0.5pt}\end{center}

\section{Appendix A. Psychology Principles and Claim-Governance
Summary}\label{appendix-a.-psychology-principles-and-claim-governance-summary}

\subsection{A.1 Domain authority}\label{a.1-domain-authority}

The twelve principles below are the controlling CU formulations for
Psychology. Domain authority identifies the formulation to cite; it does
not make conditional, clinical, or empirical content universal or final.
Consequential use remains subject to field-specific review and
affected-community participation.

In the canonical sixty-four-stage Coherence Mandala, Psychology
principally develops Stages 16--21: Representational Organization,
Anticipatory Regulation, Learning and Error-Guided Revision, Graded
Agency, Temporal Binding and Narrative, and Meaning. CU-COM15--CU-COM17
make explicit the framework commitments that follow from the meaning,
trauma, and repair architecture without turning them into clinical
mechanisms.

\subsection{CU-Ψ1 --- Meaning Space}\label{cu-ux3c81-meaning-space}

\textbf{Candidate status:} Definition / analytic construction.

For a declared psychological question, a \textbf{meaning space} is a
model-dependent set of distinctions along which events, actions,
relationships, memories, identities, and possibilities can matter to a
person or group:

\[\adjustbox{max width=\linewidth}{$\displaystyle 
\mathcal M_{\Psi}=\mathcal M_{\Psi}(\text{person},\text{context},\pi,T,\mathcal O).
$}\]

Its coordinates may include expected consequence, bodily salience,
social significance, autobiographical relevance, moral implication,
practical affordance, or sacred orientation. Its coordinates and
dimensionality must be declared; no universal dimensionality or fixed
axis set is assumed. A meaning space is an analyst's construct unless
evidence establishes that its distinctions are represented or enacted by
the person.

\textbf{Principal failure mode:} Treating the researcher's vocabulary as
the person's phenomenology.

\subsection{CU-Ψ2 --- Meaning Maps}\label{cu-ux3c82-meaning-maps}

\textbf{Candidate status:} Analytic bridge principle.

A \textbf{meaning map} is a structured, revisable relation among cues,
interpretations, affective responses, bodily states, memories,
commitments, predicted consequences, and action tendencies. Meaning maps
may be explicit or implicit, proposition-like or procedural,
individually carried or socially scaffolded. Their psychological
importance lies in how they organize attention, expectation, action, and
revision---not in their conformity to one canonical map.

\textbf{Principal failure mode:} Assuming that internal consistency,
narrative elegance, or social stability establishes descriptive truth.

\subsection{CU-Ψ3 --- Psychological Coherence
Profile}\label{cu-ux3c83-psychological-coherence-profile}

\textbf{Candidate status:} Canonical-conditional methodological
principle.

Psychological organization should ordinarily be represented by a
declared profile rather than one universal scalar. A default research
profile may include descriptive adequacy, affect tolerance and
regulation, regulatory flexibility, identity and narrative integration,
relational capacity, agency, embodied regulation, meaning or purpose,
and repair capacity:

\[\adjustbox{max width=\linewidth}{$\displaystyle 
\mathbf P_{\Psi}
=
\langle
p_{\mathrm{Ad}},p_{\mathrm{Af}},p_{\mathrm{Fr}},p_{\mathrm{In}},p_{\mathrm{Rc}},p_{\mathrm{Ga}},p_{\mathrm{Br}},p_{\mathrm{Mp}},p_{\mathrm{Qr}}
\rangle.
$}\]

The dimensions, observations, scales, and horizons must be justified for
the question. Scalarization is permitted only when its weights,
normalization, protected floors, affected persons, and decision purpose
are declared. Non-dominating profiles may remain incomparable.

\textbf{Principal failure mode:} Smuggling a moral or clinical ideal
into an apparently neutral score.

\subsection{CU-Ψ4 --- Integration Load and Regulatory
Capacity}\label{cu-ux3c84-integration-load-and-regulatory-capacity}

\textbf{Candidate status:} Model family / bridge hypothesis.

At a given time and scale, a person confronts a demand profile \(d_t\)
and possesses a context-sensitive capacity region \(\mathcal C_t\).
\textbf{Integration overload} occurs when relevant demands cannot be
metabolized, coordinated, deferred, delegated, escaped, or safely
compartmentalized within that capacity:

\[\adjustbox{max width=\linewidth}{$\displaystyle 
d_t\notin\mathcal C_t.
$}\]

Stress and distress depend not only on demand magnitude but on
predictability, controllability, meaning, bodily condition, social
support, prior learning, structural power, and temporal horizon.
Overload is not identical with disorder, and low expressed distress is
not sufficient evidence that demands have been integrated.

\textbf{Principal failure mode:} Locating in the individual a burden
produced by an unsafe or unjust environment.

\subsection{CU-Ψ5 --- Psychological Regulation and
Meaning-Seeking}\label{cu-ux3c85-psychological-regulation-and-meaning-seeking}

\textbf{Candidate status:} Empirical bridge hypothesis.

Human psychological activity is organized by multiple, sometimes
conflicting regulatory tendencies, including safety, attachment,
exploration, competence, status, identity protection, truth-seeking,
care, belonging, autonomy, pleasure, pain avoidance, and sacred or
long-horizon commitment. People often seek interpretations and practices
that make action possible, but there is no established universal
psychological drive to maximize one coherence quantity.

\textbf{Principal failure mode:} Redescribing every motive after the
fact as coherence-seeking and thereby making the principle
unfalsifiable.

\subsection{CU-Ψ6 --- Local--Global and Short--Long
Tension}\label{cu-ux3c86-localglobal-and-shortlong-tension}

\textbf{Candidate status:} Bridge principle.

A pattern may be adaptive for one subsystem, relationship, role, or
short horizon while undermining another subsystem, wider context,
affected person, or longer horizon. Psychological evaluation must
therefore specify at least:

\[\adjustbox{max width=\linewidth}{$\displaystyle 
(\text{part or whole},\ \text{context},\ \text{horizon},\ \text{protected floors}).
$}\]

Suppression, avoidance, compartmentalization, conformity, disclosure,
confrontation, and self-revision can each reverse their profile effects
as these declarations change. No local gain or immediate relief licenses
an inference to global or durable health.

\textbf{Principal failure mode:} Treating the analyst's preferred
horizon as uniquely rational while ignoring acute safety.

\subsection{CU-Ψ7 --- Repair Capacity}\label{cu-ux3c87-repair-capacity}

\textbf{Candidate status:} Bridge principle / empirical hypothesis.

\textbf{Repair capacity} is the context-sensitive ability to detect
consequential mismatch or rupture, remain or return within a tolerable
range of engagement, revise models or strategies, recruit support, make
restitution where appropriate, restore workable coordination, and learn
in ways that reduce future recurrence. Repair need not restore the prior
state. It may require grief, accommodation, boundary, identity change,
or transformation of the environment.

\textbf{Principal failure mode:} Equating repair with reconciliation,
forgiveness, symptom elimination, productivity, or return to a harmful
baseline.

\subsection{CU-Ψ8 --- Protective Compartmentalization and
Fragmentation}\label{cu-ux3c88-protective-compartmentalization-and-fragmentation}

\textbf{Candidate status:} Conditional bridge principle.

Psychological separation among memories, affects, roles, self-states, or
action systems can protect functioning when full integration would
exceed present capacity or threaten safety. The same separation can
become costly when it blocks reality contact, agency, flexible access,
relationship, or repair after the original conditions change. The
distinction between \textbf{protection} and \textbf{fragmentation} is
therefore functional, contextual, and horizon-relative rather than
purely structural.

\textbf{Principal failure mode:} Pathologizing survival adaptations or
romanticizing costly dissociation.

\subsection{CU-Ψ9 --- Developmental Path Dependence and
Plasticity}\label{cu-ux3c89-developmental-path-dependence-and-plasticity}

\textbf{Candidate status:} Empirical bridge principle.

Psychological development is path-dependent: earlier relationships,
bodily states, learning histories, institutions, and material conditions
alter the accessibility, cost, and meaning of later transitions.
Development nevertheless remains plastic. Similar adult patterns can
arise by different paths, and different outcomes can arise from similar
beginnings. Attractor language is legitimate only when a model specifies
states, transitions, perturbations, and evidence of return or basin
structure.

\textbf{Principal failure mode:} Turning statistical regularities or
culturally preferred biographies into destiny.

\subsection{CU-Ψ10 --- Attachment and
Co-Regulation}\label{cu-ux3c810-attachment-and-co-regulation}

\textbf{Candidate status:} Empirical bridge principle.

Human regulation is partly distributed across relationships. Attachment
figures and other trusted partners can alter perceived threat,
attention, effort, affect tolerance, exploration, and recovery. Repeated
co-regulation can scaffold capacities later enacted more independently,
but healthy adulthood is not independence from others. Attachment
categories are probabilistic summaries, not fixed essences or exhaustive
identities.

\textbf{Principal failure mode:} Using attachment labels as
retrospective total explanations or as grounds for blame.

\subsection{CU-Ψ11 --- Self-Deception as Motivated Meaning-Map
Distortion}\label{cu-ux3c811-self-deception-as-motivated-meaning-map-distortion}

\textbf{Candidate status:} Bridge principle.

Self-deception is a pattern of motivated, asymmetric information
processing in which a person protects an identity, attachment,
advantage, policy, or other stake by selectively excluding, reframing,
discounting, or neutralizing relevant reasons. False belief,
uncertainty, bias, trauma adaptation, sincere disagreement, or ordinary
inconsistency is not sufficient. Attribution should rely on convergent
process evidence, plausible motive, counterfactual asymmetry, and
serious comparison with alternatives.

\textbf{Principal failure mode:} Converting disagreement with the
analyst into diagnosis of corruption.

\subsection{CU-Ψ12 --- Psychological Coherence
Infrastructure}\label{cu-ux3c812-psychological-coherence-infrastructure}

\textbf{Candidate status:} Social-psychological bridge hypothesis.

Psychological capacities are materially and socially scaffolded. Stable
housing, bodily safety, time, income, accessible care, trustworthy
information, reciprocal relationships, cultural intelligibility, legal
standing, and opportunities for meaningful participation can expand or
contract what a person can regulate, revise, and repair. Population
distress should therefore be analyzed through both individual mechanisms
and the infrastructure that distributes exposure, burden, support, and
agency.

\textbf{Principal failure mode:} Using social explanation to erase
individual differences, biology, or the value of person-level treatment.

\subsection{A.14 Claim-governance
summary}\label{a.14-claim-governance-summary}

Applications must declare the person or target, context, boundary,
horizon, profile dimensions, protected floors, observation channels,
power relations, alternative explanations, clinical scope, and
intervention effects on agency. Self-deception requires convergent
process evidence. Optimization language must distinguish designer
objectives, as-if analysis, and evolutionary selection. Spiritual
interpretation may not override safety, consent, evidence, or treatment
standards.

\begin{center}\rule{0.5\linewidth}{0.5pt}\end{center}

\subsection{A.15 Commitment
Cross-Reference}\label{a.15-commitment-cross-reference}

CU-COM10 is the official Psychology domain commitment anchor for
Sections 13--14. CU-COM14--CU-COM17 and CU-COM19 provide the related
commitments to love, participatory meaning, meaning-revision, redemptive
possibility, and present enactment of Omega. None is a clinical
mechanism or treatment recommendation. None sanctifies harm, guarantees
growth, or overrides safety, consent, affected-person authority,
evidence-based care, material conditions, or substantive community
review.

\section{Appendix B. Major Revisions from
v0.09}\label{appendix-b.-major-revisions-from-v0.09}

This appendix records only the changes needed to orient readers. The
complete source manuscript, patch, independent audit, audit response,
claim ledger, and reproducibility archive are preserved in the technical
package.

\begin{longtable}[]{@{}
  >{\raggedright\arraybackslash}p{(\columnwidth - 6\tabcolsep) * \real{0.2500}}
  >{\raggedright\arraybackslash}p{(\columnwidth - 6\tabcolsep) * \real{0.2500}}
  >{\raggedright\arraybackslash}p{(\columnwidth - 6\tabcolsep) * \real{0.2500}}
  >{\raggedright\arraybackslash}p{(\columnwidth - 6\tabcolsep) * \real{0.2500}}@{}}
\toprule\noalign{}
\begin{minipage}[b]{\linewidth}\raggedright
Topic
\end{minipage} & \begin{minipage}[b]{\linewidth}\raggedright
v0.09 formulation
\end{minipage} & \begin{minipage}[b]{\linewidth}\raggedright
Version 0.10 formulation
\end{minipage} & \begin{minipage}[b]{\linewidth}\raggedright
Disposition
\end{minipage} \\
\midrule\noalign{}
\endhead
\bottomrule\noalign{}
\endlastfoot
Psychological measure & Universal Field Coherence Score &
Multidimensional profile; scalarization is purpose-relative & Universal
scalar retired \\
Stress & Inverse Entropic Stress & Integration load and regulatory
capacity under declared context & Entropy identity retired \\
Motivation & One coherence drive & Plural and sometimes conflicting
motives & Universal drive retired \\
Meaning & One coherence field or amount & Relational meaning maps
spanning description, embodiment, narrative, and commitment &
Reconstructed \\
Trauma & Universal coherence catastrophe & Heterogeneous protective and
injurious trajectories & Universal trajectory retired \\
Neurodiversity & Alternative coherence & Difference, disability,
distress, mismatch, stigma, and support assessed separately &
Reconstructed \\
Psychopathology & Persistent coherence failure & Heterogeneous
biomedical, developmental, learning, network, social, and
phenomenological models & Essence claim retired \\
Therapy & One coherence-modulation mechanism & Plural mechanisms, common
factors, person-specific formulation, and material conditions &
Reconstructed \\
Development & One hierarchy ending in integrated truth and compassion &
Multidimensional, asynchronous, culturally situated development &
Universal hierarchy retired \\
Self-deception & Broadly inferred from inconsistency or defense &
Motivated asymmetric processing requiring convergent evidence &
Narrowed \\
Spiritual meaning & Sometimes entangled with explanatory claims &
Designated interpretation and commitment sections & Register
clarified \\
\end{longtable}

\begin{center}\rule{0.5\linewidth}{0.5pt}\end{center}

\section{Appendix C. Technical Supplement and
Reproducibility}\label{appendix-c.-technical-supplement-and-reproducibility}

\textbf{Released evidence package.} The current supplement is
\href{../supplements/06_Psychology_Technical_Supplement_v0.10.md}{06\_Psychology\_Technical\_Supplement\_v0.10.md}.
Its companion evidence directory is
\href{../evidence/06_Psychology/README.md}{06\_Psychology}. The
supplement and evidence directory expose the derivations, deterministic
constructions, scripts, machine-readable results, source ledger, and
audit materials referenced in this manuscript. They do not elevate a toy
model, simulation, or reconstructed check into target-domain empirical
validation.

The accompanying technical supplement contains:

\begin{itemize}
\tightlist
\item
  formal propositions and counterexamples;
\item
  the complete claim and evidence ledger;
\item
  the twelve-principle registry and fifteen governance rules;
\item
  the deterministic horizon-reversal model;
\item
  the exact probe and boundary conventions;
\item
  coefficient, ablation, and initial-state sensitivity checks;
\item
  the independent audit and audit response;
\item
  the literature and dependency ledgers;
\item
  the full revision history and source-to-publication patch;
\item
  the reproducibility script, generated JSON, environment declaration,
  and package hashes.
\end{itemize}

The illustrative model was deliberately constructed so that immediate
suppression wins under short-horizon comfort weights and supported
engagement wins under longer-horizon functioning, capacity, and
probe-response weights. Its internal robustness does not establish
treatment efficacy or clinical generality.

\begin{center}\rule{0.5\linewidth}{0.5pt}\end{center}

\section{References}\label{references}

Ainsworth, M. D. S., Blehar, M. C., Waters, E., \& Wall, S. (1978).
\emph{Patterns of attachment: A psychological study of the strange
situation}. Erlbaum.

Alegría, M., NeMoyer, A., Falgàs Bagué, I., Wang, Y., \& Alvarez, K.
(2018). Social determinants of mental health: Where we are and where we
need to go. \emph{Current Psychiatry Reports, 20}, 95.

Bartlett, F. C. (1932). \emph{Remembering: A study in experimental and
social psychology}. Cambridge University Press.

Billings, J., \& Nicholls, H. (2025). PTSD and complex PTSD, current
treatments and debates: A review of reviews. \emph{British Medical
Bulletin, 156}(1), ldaf015. \url{https://doi.org/10.1093/bmb/ldaf015}

Bonanno, G. A., \& Burton, C. L. (2013). Regulatory flexibility: An
individual-differences perspective on coping and emotion regulation.
\emph{Perspectives on Psychological Science, 8}(6), 591--612.

Borsboom, D. (2017). A network theory of mental disorders. \emph{World
Psychiatry, 16}(1), 5--13.

Bowlby, J. (1969). \emph{Attachment and loss: Vol. 1. Attachment}. Basic
Books.

Campbell, J. D., Trapnell, P. D., Heine, S. J., Katz, I. M., Lavallee,
L. F., \& Lehman, D. R. (1996). Self-concept clarity: Measurement,
personality correlates, and cultural boundaries. \emph{Journal of
Personality and Social Psychology, 70}(1), 141--156.

Candia-Rivera, D., Engelen, T., Babo-Rebelo, M., \& Salamone, P. C.
(2024). Interoception, network physiology and the emergence of bodily
self-awareness. \emph{Neuroscience \& Biobehavioral Reviews, 165},
105864. \url{https://doi.org/10.1016/j.neubiorev.2024.105864}

Carhart-Harris, R. L., \& Friston, K. J. (2019). REBUS and the anarchic
brain: Toward a unified model of the brain action of psychedelics.
\emph{Pharmacological Reviews, 71}(3), 316--344.

Clark, A. (2013). Whatever next? Predictive brains, situated agents, and
the future of cognitive science. \emph{Behavioral and Brain Sciences,
36}(3), 181--204.

Clamor, A., Lincoln, T. M., \& Schulze, L. (2026). Emotion regulation in
mental disorders: A systematic review and multilevel meta-analysis of
transdiagnostic and disorder-specific impairments. \emph{Psychological
Bulletin, 152}(2), 157--200. \url{https://doi.org/10.1037/bul0000512}

Clemente, R., Murphy, A., \& Murphy, J. (2024). The relationship between
self-reported interoception and anxiety: A systematic review and
meta-analysis. \emph{Neuroscience \& Biobehavioral Reviews, 167},
105923. \url{https://doi.org/10.1016/j.neubiorev.2024.105923}

Coan, J. A., \& Sbarra, D. A. (2015). Social baseline theory: The social
regulation of risk and effort. \emph{Current Opinion in Psychology, 1},
87--91.

Crompton, C. J., Ropar, D., Evans-Williams, C. V. M., Flynn, E. G., \&
Fletcher-Watson, S. (2020). Autistic peer-to-peer information transfer
is highly effective. \emph{Autism, 24}(7), 1704--1712.
\url{https://doi.org/10.1177/1362361320919286}

Crompton, C. J., Foster, S. J., Wilks, C. E. H., Dodd, M., Efthimiou, T.
N., Ropar, D., Sasson, N. J., Lages, M., \& Fletcher-Watson, S. (2025).
Information transfer within and between autistic and non-autistic
people: A registered report. \emph{Nature Human Behaviour, 9}(7),
1488--1500. \url{https://doi.org/10.1038/s41562-025-02163-z}

Cuijpers, P., Miguel, C., Ciharova, M., Harrer, M., Basic, D., Cristea,
I. A., et al.~(2024). Absolute and relative outcomes of psychotherapies
for eight mental disorders: A systematic review and meta-analysis.
\emph{World Psychiatry, 23}(2), 267--275.
\url{https://doi.org/10.1002/wps.21203}

Deci, E. L., \& Ryan, R. M. (2000). The ``what'' and ``why'' of goal
pursuits: Human needs and the self-determination of behavior.
\emph{Psychological Inquiry, 11}(4), 227--268.

Elsey, J. W. B., Van Ast, V. A., \& Kindt, M. (2018). Human memory
reconsolidation: A guiding framework and critical review of the
evidence. \emph{Psychological Bulletin, 144}(8), 797--848.

Festinger, L. (1957). \emph{A theory of cognitive dissonance}. Stanford
University Press.

Fisher, A. J., Medaglia, J. D., \& Jeronimus, B. F. (2018). Lack of
group-to-individual generalizability is a threat to human subjects
research. \emph{Proceedings of the National Academy of Sciences,
115}(27), E6106--E6115.

Flückiger, C., Del Re, A. C., Wampold, B. E., \& Horvath, A. O. (2018).
The alliance in adult psychotherapy: A meta-analytic synthesis.
\emph{Psychotherapy, 55}(4), 316--340.

Friston, K. (2010). The free-energy principle: A unified brain theory?
\emph{Nature Reviews Neuroscience, 11}, 127--138.

Galatzer-Levy, I. R., Huang, S. H., \& Bonanno, G. A. (2018).
Trajectories of resilience and dysfunction following potential trauma: A
review and statistical evaluation. \emph{Clinical Psychology Review,
63}, 41--55.

Greig, L., Coundouris, S. P., \& Henry, J. D. (2026). Autistic traits
and camouflaging: A meta-analysis. \emph{Autism, 30}(6), 1398--1415.
\url{https://doi.org/10.1177/13623613261437500}

Gross, J. J. (2002). Emotion regulation: Affective, cognitive, and
social consequences. \emph{Psychophysiology, 39}(3), 281--291.

Gross, J. J. (2015). Emotion regulation: Current status and future
prospects. \emph{Psychological Inquiry, 26}(1), 1--26.

Hofmann, S. G., \& Hayes, S. C. (2019). The future of intervention
science: Process-based therapy. \emph{Clinical Psychological Science,
7}(1), 37--50.

Hohwy, J. (2013). \emph{The predictive mind}. Oxford University Press.

Kashdan, T. B., \& Rottenberg, J. (2010). Psychological flexibility as a
fundamental aspect of health. \emph{Clinical Psychology Review, 30}(7),
865--878.

King, L. A., \& Hicks, J. A. (2021). The science of meaning in life.
\emph{Annual Review of Psychology, 72}, 561--584.

Kotov, R., Krueger, R. F., Watson, D., Achenbach, T. M., Althoff, R. R.,
Bagby, R. M., et al.~(2017). The Hierarchical Taxonomy of
Psychopathology (HiTOP): A dimensional alternative to traditional
nosologies. \emph{Journal of Abnormal Psychology, 126}(4), 454--477.

Lee, J. L. C., Nader, K., \& Schiller, D. (2017). An update on memory
reconsolidation updating. \emph{Trends in Cognitive Sciences, 21}(7),
531--545.

Lind, M., Ture, S., McAdams, D. P., \& Cowan, H. R. (2024). Narrative
identity, traits, and trajectories of depression and well-being: A
9-year longitudinal study. \emph{Psychological Science, 35}(12),
1325--1339. \url{https://doi.org/10.1177/09567976241296512}

Lindahl, J. R., Fisher, N. E., Cooper, D. J., Rosen, R. K., \& Britton,
W. B. (2017). The varieties of contemplative experience: A mixed-methods
study of meditation-related challenges in Western Buddhists. \emph{PLOS
ONE, 12}(5), e0176239.

Loftus, E. F. (2005). Planting misinformation in the human mind: A
30-year investigation of the malleability of memory. \emph{Learning \&
Memory, 12}(4), 361--366.

Lund, C., Brooke-Sumner, C., Baingana, F., Baron, E. C., Breuer, E.,
Chandra, P., et al.~(2018). Social determinants of mental disorders and
the Sustainable Development Goals: A systematic review of reviews.
\emph{The Lancet Psychiatry, 5}(4), 357--369.

Lunansky, G., Bonanno, G. A., Blanken, T. F., van Borkulo, C. D.,
Cramer, A. O. J., \& Borsboom, D. (2025). Bouncing back from life's
perturbations: Formalizing psychological resilience from a complex
systems perspective. \emph{Psychological Review, 132}(6), 1396--1409.
\url{https://doi.org/10.1037/rev0000497}

Macri, J. A., \& Rogge, R. D. (2024). Examining domains of psychological
flexibility and inflexibility as treatment mechanisms in acceptance and
commitment therapy: A comprehensive systematic and meta-analytic review.
\emph{Clinical Psychology Review, 110}, 102432.
\url{https://doi.org/10.1016/j.cpr.2024.102432}

Martela, F., \& Steger, M. F. (2016). The three meanings of meaning in
life: Distinguishing coherence, purpose, and significance. \emph{The
Journal of Positive Psychology, 11}(5), 531--545.

McAdams, D. P. (2001). The psychology of life stories. \emph{Review of
General Psychology, 5}(2), 100--122.

McAdams, D. P., \& McLean, K. C. (2013). Narrative identity.
\emph{Current Directions in Psychological Science, 22}(3), 233--238.

Nel, J., Spedding, M., \& Malcolm-Smith, S. (2025). Consolidating a
framework of autistic camouflaging strategies: An integrative systematic
review. \emph{Autism, 29}(10), 2379--2394.
\url{https://doi.org/10.1177/13623613251335472}

Nejadghaderi, S. A., Mousavi, S. E., Fazlollahi, A., Motlagh Asghari,
K., \& Garfin, D. R. (2024). Efficacy of yoga for posttraumatic stress
disorder: A systematic review and meta-analysis of randomized controlled
trials. \emph{Psychiatry Research, 340}, 116098.
\url{https://doi.org/10.1016/j.psychres.2024.116098}

Pellicano, E., \& den Houting, J. (2022). Annual research review:
Shifting from ``normal science'' to neurodiversity in autism science.
\emph{Journal of Child Psychology and Psychiatry, 63}(4), 381--396.

Rader, G. K. D. (2026). \emph{Coherence Universalism:
Physics---Coherence, Accessibility, and the Ordered Emergence of
Physical Description} (Version 0.10). Heaven≡Earth Foundation.

Rader, G. K. D. (2026). \emph{Coherence Universalism: Metaphysics and
Epistemology---Coherence Under Constraint, Explanatory Order, and
Warrant Across Registers} (Version 0.10). Heaven≡Earth Foundation.

Rader, G. K. D. (2026). \emph{Coherence Universalism:
Foundations---Principles, Domain Gateways, Commitments, and Constraints
on Responsible Application} (Version 0.10). Heaven≡Earth Foundation.

Rader, G. K. D. (2026). \emph{Coherence Universalism: Biology---Life as
Active Self-Maintenance: Viability, Individuation, and Multiscale
Regulation} (Version 0.10). Heaven≡Earth Foundation.

Rader, G. K. D. (2026). \emph{Coherence Universalism:
Consciousness---Perspective, Temporal Depth, Valence, and the
Architecture of Experience} (Version 0.10). Heaven≡Earth Foundation.

Rader, G. K. D. (2026). \emph{Denying Knowledge to Make Room for
Commitment---Truth, Commitment, and Nonmonotonic Warrant} (Version
0.10). Heaven≡Earth Foundation.

Salkovskis, P. M., Sighvatsson, M. B., \& Sigurdsson, J. F. (2024). How
effective psychological treatments work: Mechanisms of change in
cognitive behavioural therapy and beyond. \emph{Behavioural and
Cognitive Psychotherapy, 51}(6), 595--615.
\url{https://doi.org/10.1017/S1352465823000590}

Sbarra, D. A., \& Hazan, C. (2008). Coregulation, dysregulation,
self-regulation: An integrative analysis and empirical agenda for
understanding adult attachment, separation, loss, and recovery.
\emph{Personality and Social Psychology Review, 12}(2), 141--167.

Schacter, D. L. (1996). \emph{Searching for memory: The brain, the mind,
and the past}. Basic Books.

Schaug, J. P., Møller, L., Reinholt, N., Illum, D. B., Græbe, F. L.,
Mikkelsen, L. B., et al.~(2025). Psychotherapies for adults with complex
presentations of PTSD: A clinical guideline and five systematic reviews
with meta-analyses. \emph{BMJ Mental Health, 28}(1), e301158.
\url{https://doi.org/10.1136/bmjment-2024-301158}

Schumacher, L., Burger, J., Echterhoff, J., \& Kriston, L. (2024).
Methodological and statistical practices of using symptom networks to
evaluate mental health interventions: A review and reflections.
\emph{Multivariate Behavioral Research, 59}(4), 663--676.
\url{https://doi.org/10.1080/00273171.2024.2335401}

Steger, M. F., Frazier, P., Oishi, S., \& Kaler, M. (2006). The Meaning
in Life Questionnaire: Assessing the presence of and search for meaning
in life. \emph{Journal of Counseling Psychology, 53}(1), 80--93.

U.S. Department of Veterans Affairs \& U.S. Department of Defense.
(2023). \emph{VA/DoD clinical practice guideline for management of
posttraumatic stress disorder and acute stress disorder}. Authors.

Wampold, B. E., \& Imel, Z. E. (2015). \emph{The great psychotherapy
debate: The evidence for what makes psychotherapy work} (2nd ed.).
Routledge.

Westhoff, M., Berg, M., Reif, A., Rief, W., \& Hofmann, S. G. (2024).
Major problems in clinical psychological science and how to address
them: Introducing a multimodal dynamical network approach.
\emph{Cognitive Therapy and Research, 48}(5), 791--807.
\url{https://doi.org/10.1007/s10608-024-10487-9}

World Health Organization. (2022). \emph{World mental health report:
Transforming mental health for all}. Author.

World Health Organization. (2024). \emph{Clinical descriptions and
diagnostic requirements for ICD-11 mental, behavioural and
neurodevelopmental disorders}. Author.

\begin{center}\rule{0.5\linewidth}{0.5pt}\end{center}

\end{document}
